%% Generated by Sphinx.
\def\sphinxdocclass{report}
\documentclass[letterpaper,10pt,english]{sphinxmanual}
\ifdefined\pdfpxdimen
   \let\sphinxpxdimen\pdfpxdimen\else\newdimen\sphinxpxdimen
\fi \sphinxpxdimen=.75bp\relax
\ifdefined\pdfimageresolution
    \pdfimageresolution= \numexpr \dimexpr1in\relax/\sphinxpxdimen\relax
\fi
%% let collapsible pdf bookmarks panel have high depth per default
\PassOptionsToPackage{bookmarksdepth=5}{hyperref}

\PassOptionsToPackage{booktabs}{sphinx}
\PassOptionsToPackage{colorrows}{sphinx}

\PassOptionsToPackage{warn}{textcomp}
\usepackage[utf8]{inputenc}
\ifdefined\DeclareUnicodeCharacter
% support both utf8 and utf8x syntaxes
  \ifdefined\DeclareUnicodeCharacterAsOptional
    \def\sphinxDUC#1{\DeclareUnicodeCharacter{"#1}}
  \else
    \let\sphinxDUC\DeclareUnicodeCharacter
  \fi
  \sphinxDUC{00A0}{\nobreakspace}
  \sphinxDUC{2500}{\sphinxunichar{2500}}
  \sphinxDUC{2502}{\sphinxunichar{2502}}
  \sphinxDUC{2514}{\sphinxunichar{2514}}
  \sphinxDUC{251C}{\sphinxunichar{251C}}
  \sphinxDUC{2572}{\textbackslash}
\fi
\usepackage{cmap}
\usepackage[T1]{fontenc}
\usepackage{amsmath,amssymb,amstext}
\usepackage{babel}



\usepackage{tgtermes}
\usepackage{tgheros}
\renewcommand{\ttdefault}{txtt}



\usepackage[Bjarne]{fncychap}
\usepackage{sphinx}

\fvset{fontsize=auto}
\usepackage{geometry}


% Include hyperref last.
\usepackage{hyperref}
% Fix anchor placement for figures with captions.
\usepackage{hypcap}% it must be loaded after hyperref.
% Set up styles of URL: it should be placed after hyperref.
\urlstyle{same}

\addto\captionsenglish{\renewcommand{\contentsname}{Contents:}}

\usepackage{sphinxmessages}
\setcounter{tocdepth}{1}



\title{WLCF}
\date{May 06, 2026}
\release{1.0.0}
\author{Alejandro Avilés}
\newcommand{\sphinxlogo}{\vbox{}}
\renewcommand{\releasename}{Release}
\makeindex
\begin{document}

\ifdefined\shorthandoff
  \ifnum\catcode`\=\string=\active\shorthandoff{=}\fi
  \ifnum\catcode`\"=\active\shorthandoff{"}\fi
\fi

\pagestyle{empty}
\sphinxmaketitle
\pagestyle{plain}
\sphinxtableofcontents
\pagestyle{normal}
\phantomsection\label{\detokenize{index::doc}}


\sphinxstepscope


\chapter{Overview}
\label{\detokenize{overview:overview}}\label{\detokenize{overview::doc}}

\section{wlcf: Weak Lensing Correlation Function}
\label{\detokenize{overview:wlcf-weak-lensing-correlation-function}}
\sphinxAtStartPar
\sphinxstylestrong{Authors:} A. Aviles, J.C. Hidalgo, E.A. Moreno\sphinxhyphen{}Alcala, G. Niz, S. Ramirez, Mario A. Rodriguez\sphinxhyphen{}Meza, S. Samario\sphinxhyphen{}Nava

\sphinxAtStartPar
For download and information, see \sphinxurl{https://github.com/rodriguezmeza/wlcf.git} .


\section{Introduction}
\label{\detokenize{overview:introduction}}
\sphinxAtStartPar
\sphinxstylestrong{wlcf} is a C code designed to compute cosmological correlation functions using theoretical frameworks such as Standard Perturbation Theory (SPT), Effective Field Theory (EFT), and the Halo Model. Currently, it focuses on the computation of the \sphinxstylestrong{three\sphinxhyphen{}point correlation function (3PCF)} for galaxy weak lensing convergence in a plane\sphinxhyphen{}wave expansion.
Paper available at (\sphinxurl{https://arxiv.org/abs/2408.16847}).


\section{Compiling and getting started}
\label{\detokenize{overview:compiling-and-getting-started}}
\sphinxAtStartPar
Download the code by cloning it from:
\sphinxurl{https://github.com/rodriguezmeza/wlcf.git} .

\sphinxAtStartPar
\sphinxstylestrong{Dependencies:}
\sphinxtitleref{wlcf} optionally requires \sphinxstylestrong{GSL (GNU Scientific Library)} and \sphinxstylestrong{FFTW3} installed in your system. Make sure to adapt the paths to these libraries in \sphinxtitleref{Makefile\_machine}.

\sphinxAtStartPar
Tu run the code go to the \sphinxtitleref{wlcf} directory and compile the code:

\begin{sphinxVerbatim}[commandchars=\\\{\}]
\PYGZdl{} cd wlcf
\PYGZdl{} make clean
\PYGZdl{} make all
\end{sphinxVerbatim}

\sphinxAtStartPar
The code has been tested with modern versions of \sphinxtitleref{gcc}. OpenMP support is optional but recommended for parallel execution. In addition, the code provides a Python wrapper that allows it to be run and controlled directly from Python.

\sphinxAtStartPar
To verify that the code runs correctly:

\begin{sphinxVerbatim}[commandchars=\\\{\}]
\PYGZdl{} cd tests
\PYGZdl{} ./wlcf
\end{sphinxVerbatim}

\sphinxAtStartPar
This will execute the code using default parameters. A directory named \sphinxtitleref{Output} will be created inside \sphinxtitleref{tests}, where all results are stored, including:
\begin{itemize}
\item {} 
\sphinxAtStartPar
output data files (e.g., \sphinxtitleref{zetam*}, \sphinxtitleref{Bmells*}, \sphinxtitleref{Bnk*})

\item {} 
\sphinxAtStartPar
a log file inside \sphinxtitleref{Output/tmp}

\end{itemize}

\sphinxAtStartPar
You can consult man page:

\begin{sphinxVerbatim}[commandchars=\\\{\}]
\PYGZdl{} man ../docs/wlcf.1
\end{sphinxVerbatim}

\sphinxAtStartPar
or open with a browser the html file: docs/man/wlcf.html


\section{License}
\label{\detokenize{overview:license}}
\sphinxAtStartPar
wlcf is written by A. Aviles et al., and is distributed under the \sphinxhref{https://github.com/rodriguezmeza/wlcf/blob/main/LICENSE}{MIT license}. If you use this program in research work that results in publications, please cite the following paper:

\sphinxAtStartPar
\sphinxhref{https://arxiv.org/abs/2408.16847}{Abraham Arvizu et al., arXiv:2048.16847}


\section{Acknowledgements}
\label{\detokenize{overview:acknowledgements}}
\sphinxAtStartPar
wlcf use/is based on the following codes or projects:
\begin{itemize}
\item {} 
\sphinxAtStartPar
\sphinxhref{https://github.com/xfangcosmo/2DFFTLog}{FFTLog}

\item {} 
\sphinxAtStartPar
\sphinxhref{https://cosmo.phys.hirosaki-u.ac.jp/takahasi/codes\_e.htm}{The BiHaloFit model of Takahashi}

\item {} 
\sphinxAtStartPar
\sphinxhref{https://home.ifa.hawaii.edu/users/barnes/zeno/index.html}{Zeno}

\item {} 
\sphinxAtStartPar
\sphinxhref{https://numerical.recipes/}{Numerical recipies}

\item {} 
\sphinxAtStartPar
\sphinxhref{https://www.gnu.org/software/gsl/}{GSL}

\item {} 
\sphinxAtStartPar
\sphinxhref{https://github.com/lesgourg/class\_public}{CLASS}

\end{itemize}

\sphinxstepscope


\chapter{WLCF Installation}
\label{\detokenize{installation:wlcf-installation}}\label{\detokenize{installation::doc}}
\sphinxAtStartPar
The wlcf code is a C\sphinxhyphen{}based implementation for computing higher\sphinxhyphen{}order correlation functions in cosmology, with optional support for parallel execution via OpenMP and a Python interface for flexible workflows.
This section guides you through installing the code, compiling it, and running a first test example.


\section{Requirements}
\label{\detokenize{installation:requirements}}
\sphinxAtStartPar
TThe code requires:
\begin{itemize}
\item {} 
\sphinxAtStartPar
A C compiler (recommended: \sphinxtitleref{gcc \textgreater{}= 10})

\item {} 
\sphinxAtStartPar
Optional libraries:
\begin{itemize}
\item {} 
\sphinxAtStartPar
\sphinxstylestrong{GSL} (GNU Scientific Library)

\item {} 
\sphinxAtStartPar
\sphinxstylestrong{FFTW3}

\end{itemize}

\item {} 
\sphinxAtStartPar
Python (optional, for the wrapper interface):
\begin{itemize}
\item {} 
\sphinxAtStartPar
\sphinxtitleref{numpy}

\item {} 
\sphinxAtStartPar
\sphinxtitleref{cython}

\end{itemize}

\end{itemize}


\section{Compilation}
\label{\detokenize{installation:compilation}}
\sphinxAtStartPar
Clone the repository:

\begin{sphinxVerbatim}[commandchars=\\\{\}]
git\PYG{+w}{ }clone\PYG{+w}{ }https://github.com/your\PYGZhy{}repo/wlcf.git
\PYG{n+nb}{cd}\PYG{+w}{ }wlcf
\end{sphinxVerbatim}

\sphinxAtStartPar
Edit the file \sphinxcode{\sphinxupquote{Makefile\_machine}} if needed to specify library paths (GSL, FFTW).

\sphinxAtStartPar
Compile the code:

\begin{sphinxVerbatim}[commandchars=\\\{\}]
make\PYG{+w}{ }clean
make
\end{sphinxVerbatim}

\sphinxAtStartPar
If GSL is not available, it can be disabled in \sphinxtitleref{Makefile\_settings} by setting:

\begin{sphinxVerbatim}[commandchars=\\\{\}]
\PYG{n+nv}{USEGSL}\PYG{+w}{ }\PYG{o}{=}\PYG{+w}{ }\PYG{l+m}{0}
\end{sphinxVerbatim}

\sphinxAtStartPar
If compilation is successful, the executable \sphinxcode{\sphinxupquote{wlcf}} will be created in the main directory.

\sphinxAtStartPar
To verify the installation:

\begin{sphinxVerbatim}[commandchars=\\\{\}]
\PYG{n+nb}{cd}\PYG{+w}{ }tests
./wlcf
\end{sphinxVerbatim}

\sphinxAtStartPar
This runs the code using default parameters.

\sphinxAtStartPar
After execution, the following will be generated:
\begin{itemize}
\item {} 
\sphinxAtStartPar
\sphinxtitleref{Output/} directory

\item {} 
\sphinxAtStartPar
Correlation function files:
\begin{itemize}
\item {} 
\sphinxAtStartPar
\sphinxtitleref{zetam*}

\item {} 
\sphinxAtStartPar
\sphinxtitleref{Bmells*}

\item {} 
\sphinxAtStartPar
\sphinxtitleref{Bnk*}

\end{itemize}

\item {} 
\sphinxAtStartPar
Log files (inside \sphinxtitleref{Output/tmp})

\item {} 
\sphinxAtStartPar
Parameter file: \sphinxtitleref{parameters\_null\sphinxhyphen{}cballs\sphinxhyphen{}usedvalues*}

\end{itemize}

\sphinxAtStartPar
This parameter file can be used as a template for custom runs.

\sphinxAtStartPar
\sphinxstylestrong{Python wrapper}

\sphinxAtStartPar
The code includes a Python interface for easier control and integration.

\sphinxAtStartPar
Build the wrapper

\begin{sphinxVerbatim}[commandchars=\\\{\}]
make\PYG{+w}{ }all
\end{sphinxVerbatim}

\sphinxAtStartPar
Install dependencies if needed:

\sphinxAtStartPar
You can use the wrapper as

\begin{sphinxVerbatim}[commandchars=\\\{\}]
\PYG{k+kn}{from}\PYG{+w}{ }\PYG{n+nn}{wlcfpy}\PYG{+w}{ }\PYG{k+kn}{import} \PYG{n}{wlcf}

\PYG{n}{w} \PYG{o}{=} \PYG{n}{wlcf}\PYG{p}{(}\PYG{p}{)}
\PYG{n}{w}\PYG{o}{.}\PYG{n}{set}\PYG{p}{(}\PYG{n}{numberThreads}\PYG{o}{=}\PYG{l+m+mi}{4}\PYG{p}{,} \PYG{n}{verbose}\PYG{o}{=}\PYG{l+m+mi}{1}\PYG{p}{)}
\PYG{n}{w}\PYG{o}{.}\PYG{n}{Run}\PYG{p}{(}\PYG{p}{)}
\end{sphinxVerbatim}

\sphinxAtStartPar
\sphinxstylestrong{Parallel execution (OpenMP)}

\sphinxAtStartPar
The code supports shared\sphinxhyphen{}memory parallelism using OpenMP.

\sphinxAtStartPar
To control the number of threads:

\begin{sphinxVerbatim}[commandchars=\\\{\}]
\PYG{n+nb}{export}\PYG{+w}{ }\PYG{n+nv}{OMP\PYGZus{}NUM\PYGZus{}THREADS}\PYG{o}{=}\PYG{l+m}{8}
\end{sphinxVerbatim}

\sphinxAtStartPar
The code may also set the number of threads internally using the parameter \sphinxtitleref{numberThreads}.

\sphinxAtStartPar
\sphinxstylestrong{HPC example (Perlmutter)}

\sphinxAtStartPar
On \sphinxstylestrong{NERSC Perlmutter}, a typical workflow in an interactive node is:

\begin{sphinxVerbatim}[commandchars=\\\{\}]
module\PYG{+w}{ }load\PYG{+w}{ }gcc
salloc\PYG{+w}{ }\PYGZhy{}N\PYG{+w}{ }\PYG{l+m}{1}\PYG{+w}{ }\PYGZhy{}C\PYG{+w}{ }cpu\PYG{+w}{ }\PYGZhy{}q\PYG{+w}{ }interactive\PYG{+w}{ }\PYGZhy{}t\PYG{+w}{ }\PYG{l+m}{01}:00:00
\PYG{n+nb}{export}\PYG{+w}{ }\PYG{n+nv}{OMP\PYGZus{}NUM\PYGZus{}THREADS}\PYG{o}{=}\PYG{l+m}{32}
\PYG{n+nb}{cd}\PYG{+w}{ }wlcf/tests
../wlcf
\end{sphinxVerbatim}

\sphinxAtStartPar
\sphinxstylestrong{Notes}
\begin{itemize}
\item {} 
\sphinxAtStartPar
If compilation fails, verify library paths in \sphinxtitleref{Makefile\_machine}

\item {} 
\sphinxAtStartPar
OpenMP requires a compatible compiler (e.g., \sphinxtitleref{gcc} with \sphinxtitleref{\sphinxhyphen{}fopenmp})

\item {} 
\sphinxAtStartPar
Performance may saturate at high thread counts depending on the workload

\item {} 
\sphinxAtStartPar
For large runs, always use compute nodes rather than login nodes

\end{itemize}

\sphinxstepscope


\chapter{Parameters}
\label{\detokenize{params:parameters}}\label{\detokenize{params::doc}}
\sphinxAtStartPar
This section describes the various parameters \sphinxstylestrong{wlcf} needs for controlling
the computation of the three\sphinxhyphen{}point correlation function (3PCF) for weak lensing convergence.


\section{Parameters related to the cosmological model}
\label{\detokenize{params:parameters-related-to-the-cosmological-model}}\begin{quote}\begin{description}
\sphinxlineitem{z}
\sphinxAtStartPar
(float) Redshift at which the correlation functions are evaluated.

\sphinxlineitem{h}
\sphinxAtStartPar
(float) Dimensionless Hubble parameter.

\sphinxlineitem{sigma8}
\sphinxAtStartPar
(float) Normalization of the matter power spectrum.

\sphinxlineitem{Omb}
\sphinxAtStartPar
(float) Baryon density parameter.

\sphinxlineitem{Omc}
\sphinxAtStartPar
(float) Cold dark matter density parameter.

\sphinxlineitem{Omnu}
\sphinxAtStartPar
(float, default=0.0) Massive neutrino density parameter.

\sphinxlineitem{ns}
\sphinxAtStartPar
(float) Spectral index of the primordial power spectrum.

\sphinxlineitem{w}
\sphinxAtStartPar
(float, default=\sphinxhyphen{}1.0) Dark energy equation of state parameter.

\end{description}\end{quote}


\section{Parameters to control the input power spectrum}
\label{\detokenize{params:parameters-to-control-the-input-power-spectrum}}\begin{quote}\begin{description}
\sphinxlineitem{fnamePS}
\sphinxAtStartPar
(str) {[}alias: ps{]} File containing the linear matter power spectrum ( P(k) ).

\end{description}\end{quote}

\sphinxAtStartPar
Use it as:

\begin{sphinxVerbatim}[commandchars=\\\{\}]
\PYG{n}{fnamePS}\PYG{o}{=}\PYG{o}{.}\PYG{o}{/}\PYG{n+nb}{input}\PYG{o}{/}\PYG{n}{linear\PYGZus{}pk}\PYG{o}{.}\PYG{n}{txt}
\end{sphinxVerbatim}


\section{Parameters related to the weak lensing kernel}
\label{\detokenize{params:parameters-related-to-the-weak-lensing-kernel}}\begin{quote}\begin{description}
\sphinxlineitem{Wg}
\sphinxAtStartPar
(int, default=0) Controls how the lensing kernel is defined.

\end{description}\end{quote}
\begin{itemize}
\item {} 
\sphinxAtStartPar
\sphinxcode{\sphinxupquote{0}} = Dirac delta at the redshift bin

\item {} 
\sphinxAtStartPar
\sphinxcode{\sphinxupquote{1}} = Read kernel from file

\end{itemize}
\begin{quote}\begin{description}
\sphinxlineitem{fWgchi}
\sphinxAtStartPar
(str) {[}alias: fwgchi{]} File containing the weak lensing kernel ( W\_g(chi) ).

\end{description}\end{quote}

\sphinxAtStartPar
Use it as:

\begin{sphinxVerbatim}[commandchars=\\\{\}]
\PYG{n}{fWgchi}\PYG{o}{=}\PYG{o}{.}\PYG{o}{/}\PYG{n+nb}{input}\PYG{o}{/}\PYG{n}{Wg\PYGZus{}Takahashi\PYGZus{}z05078}\PYG{o}{.}\PYG{n}{txt}
\end{sphinxVerbatim}


\section{Parameters controlling the model and numerical setup}
\label{\detokenize{params:parameters-controlling-the-model-and-numerical-setup}}\begin{quote}\begin{description}
\sphinxlineitem{tree\_level}
\sphinxAtStartPar
(int, default=3) {[}alias: tlev{]} Model used for the bispectrum calculation.

\end{description}\end{quote}
\begin{itemize}
\item {} 
\sphinxAtStartPar
\sphinxcode{\sphinxupquote{1}} = Standard Perturbation Theory (SPT)

\item {} 
\sphinxAtStartPar
\sphinxcode{\sphinxupquote{2}} = Tree\sphinxhyphen{}level approximation

\item {} 
\sphinxAtStartPar
\sphinxcode{\sphinxupquote{3}} = Effective Field Theory (EFT)

\item {} 
\sphinxAtStartPar
\sphinxcode{\sphinxupquote{4}} = Halo Model (HM)

\end{itemize}
\begin{quote}\begin{description}
\sphinxlineitem{zbin}
\sphinxAtStartPar
(float) Redshift bin used in the projection.

\sphinxlineitem{mMax}
\sphinxAtStartPar
(int) Maximum multipole order for the expansion.

\sphinxlineitem{chiQuadSteps}
\sphinxAtStartPar
(int) {[}alias: chiqst{]} Number of integration steps over comoving distance.

\sphinxlineitem{GLpoints}
\sphinxAtStartPar
(int) {[}alias: gl{]} Number of Gauss\textendash{}Legendre points.

\sphinxlineitem{Nell}
\sphinxAtStartPar
(int) Number of angular multipoles.

\sphinxlineitem{ellmin}
\sphinxAtStartPar
(float) Minimum multipole value.

\sphinxlineitem{ellmax}
\sphinxAtStartPar
(float) Maximum multipole value.

\end{description}\end{quote}


\section{Parameters to control output files}
\label{\detokenize{params:parameters-to-control-output-files}}\begin{quote}\begin{description}
\sphinxlineitem{rootDir}
\sphinxAtStartPar
(str, default=”Output”) {[}alias: root{]} Directory where output files will be written.

\sphinxlineitem{prefix}
\sphinxAtStartPar
(str, default=”{\color{red}\bfseries{}run1\_}”) {[}alias: pre{]} Prefix for output file names.

\sphinxlineitem{path\_Bells}
\sphinxAtStartPar
(str, default=”Bell\_outputs”) {[}alias: bellout{]} Directory where angular correlation functions are stored.

\sphinxlineitem{writevectors}
\sphinxAtStartPar
(bool, default=true) Whether to write intermediate vectors.

\end{description}\end{quote}


\section{Miscellaneous parameters}
\label{\detokenize{params:miscellaneous-parameters}}\begin{quote}\begin{description}
\sphinxlineitem{verbose}
\sphinxAtStartPar
(int, default=2) {[}alias: verb{]} Controls the amount of information printed to the terminal.

\end{description}\end{quote}
\begin{itemize}
\item {} 
\sphinxAtStartPar
0 = no output unless there is an error

\item {} 
\sphinxAtStartPar
1 = warnings

\item {} 
\sphinxAtStartPar
2 = progress information

\item {} 
\sphinxAtStartPar
3 = debugging output

\end{itemize}
\begin{quote}\begin{description}
\sphinxlineitem{verbose\_log}
\sphinxAtStartPar
(int, default=1) {[}alias: verblog{]} Controls logging to file inside \sphinxtitleref{Output/tmp}.

\sphinxlineitem{chatty}
\sphinxAtStartPar
(int, default=2) Additional verbosity level for internal messages.

\sphinxlineitem{numberThreads}
\sphinxAtStartPar
(int, default=4) {[}alias: nthreads{]} Number of OpenMP threads used during execution.

\end{description}\end{quote}

\sphinxAtStartPar
It is required to enable OpenMP by setting:

\begin{sphinxVerbatim}[commandchars=\\\{\}]
\PYG{n}{OPENMPMACHINE} \PYG{o}{=} \PYG{l+m+mi}{1}
\end{sphinxVerbatim}

\sphinxAtStartPar
in \sphinxcode{\sphinxupquote{Makefile\_settings}} and recompiling the code.
\begin{quote}\begin{description}
\sphinxlineitem{options}
\sphinxAtStartPar
(str, default=””) {[}alias: opt{]} Additional options controlling code behavior.

\end{description}\end{quote}
\begin{itemize}
\item {} 
\sphinxAtStartPar
It is not necessary to specify all parameters. Only the relevant ones need to be provided. The rest will take default values.

\item {} 
\sphinxAtStartPar
A parameter file is automatically generated after each run:

\begin{sphinxVerbatim}[commandchars=\\\{\}]
\PYG{n}{Output}\PYG{o}{/}\PYG{n}{parameters\PYGZus{}null}\PYG{o}{\PYGZhy{}}\PYG{n}{cballs}\PYG{o}{\PYGZhy{}}\PYG{n}{usedvalues}
\end{sphinxVerbatim}

\end{itemize}

\sphinxAtStartPar
This file can be used as a template for custom configurations.
\begin{itemize}
\item {} 
\sphinxAtStartPar
When specifying \sphinxtitleref{rootDir}, avoid using \sphinxtitleref{./} at the beginning or \sphinxtitleref{/} at the end if it is a local directory.

\end{itemize}

\sphinxstepscope


\chapter{Linear power spectrum files}
\label{\detokenize{ps_files:linear-power-spectrum-files}}\label{\detokenize{ps_files::doc}}
\sphinxAtStartPar
This section describes the format of the linear power spectrum required to compute correlation functions using \sphinxstylestrong{wlcf}.

\sphinxAtStartPar
\sphinxstylestrong{Standard formats}
\begin{quote}\begin{description}
\sphinxlineitem{columns\sphinxhyphen{}ascii}
\sphinxAtStartPar
The linear power spectrum is provided as an ASCII file with two columns:

\end{description}\end{quote}
\begin{itemize}
\item {} 
\sphinxAtStartPar
\sphinxcode{\sphinxupquote{k}} = wavenumber in units of ( h ,mathrm\{Mpc\}\textasciicircum{}\{\sphinxhyphen{}1\} )

\item {} 
\sphinxAtStartPar
\sphinxcode{\sphinxupquote{P(k)}} = linear matter power spectrum in units of ( (mathrm\{Mpc\}/h)\textasciicircum{}3 )

\end{itemize}

\sphinxAtStartPar
Each row corresponds to a value of ( k ) and its associated power spectrum.

\sphinxAtStartPar
Example:

\begin{sphinxVerbatim}[commandchars=\\\{\}]
\PYG{c+c1}{\PYGZsh{} k [h/Mpc]    P(k) [(Mpc/h)\PYGZca{}3]}
\PYG{l+m+mf}{1.000000e\PYGZhy{}03}   \PYG{l+m+mf}{2.345678e+04}
\PYG{l+m+mf}{2.000000e\PYGZhy{}03}   \PYG{l+m+mf}{1.987654e+04}
\PYG{l+m+mf}{5.000000e\PYGZhy{}03}   \PYG{l+m+mf}{1.234567e+04}
\PYG{o}{.}\PYG{o}{.}\PYG{o}{.}
\end{sphinxVerbatim}

\sphinxAtStartPar
\sphinxstylestrong{Usage in wlcf}

\sphinxAtStartPar
To use a power spectrum file, specify:

\begin{sphinxVerbatim}[commandchars=\\\{\}]
\PYG{n}{fnamePS}\PYG{o}{=}\PYG{o}{.}\PYG{o}{/}\PYG{n+nb}{input}\PYG{o}{/}\PYG{n}{linear\PYGZus{}pk}\PYG{o}{.}\PYG{n}{txt}
\end{sphinxVerbatim}

\sphinxAtStartPar
The linear power spectrum can be generated using external Boltzmann solvers such as:
\begin{itemize}
\item {} 
\sphinxAtStartPar
\sphinxstylestrong{CAMB}

\item {} 
\sphinxAtStartPar
\sphinxstylestrong{CLASS}

\end{itemize}

\sphinxAtStartPar
These tools allow you to compute ( P(k) ) for a given cosmological model.

\sphinxAtStartPar
Make sure to export the power spectrum at the desired redshift and in the correct units.

\sphinxstepscope


\chapter{Pre/Post processing}
\label{\detokenize{pre_post_process:pre-post-processing}}\label{\detokenize{pre_post_process::doc}}
\sphinxAtStartPar
This section describes the various pre/pos processing for controlling
what the \sphinxtitleref{wlcf} before/after main processing:


\section{\sphinxstylestrong{Pre\sphinxhyphen{}processing}}
\label{\detokenize{pre_post_process:pre-processing}}
\sphinxAtStartPar
Pre\sphinxhyphen{}processing involves preparing all required inputs before running \sphinxstylestrong{wlcf}.

\sphinxAtStartPar
Typical tasks include:
\begin{itemize}
\item {} 
\sphinxAtStartPar
Generating the \sphinxstylestrong{linear power spectrum} ( P(k) ) using external tools (e.g., CAMB or CLASS)

\item {} 
\sphinxAtStartPar
Preparing the \sphinxstylestrong{weak lensing kernel} file (if \sphinxtitleref{Wg=1})

\item {} 
\sphinxAtStartPar
Setting cosmological parameters consistently with the input data

\item {} 
\sphinxAtStartPar
Organizing input files into accessible directories

\end{itemize}

\sphinxAtStartPar
The accuracy of the final correlation functions strongly depends on the quality and consistency of the input power spectrum and cosmological parameters.


\section{\sphinxstylestrong{Post\sphinxhyphen{}processing}}
\label{\detokenize{pre_post_process:post-processing}}
\sphinxAtStartPar
Post\sphinxhyphen{}processing refers to analyzing and visualizing the outputs produced by \sphinxstylestrong{wlcf}.

\sphinxAtStartPar
Typical tasks include:
\begin{itemize}
\item {} 
\sphinxAtStartPar
Reading output files such as:
\begin{itemize}
\item {} 
\sphinxAtStartPar
\sphinxtitleref{zetam*} (angular correlation functions)

\item {} 
\sphinxAtStartPar
\sphinxtitleref{Bmells*} (projected bispectrum)

\item {} 
\sphinxAtStartPar
\sphinxtitleref{Bnk*} (bispectrum in Fourier space)

\end{itemize}

\item {} 
\sphinxAtStartPar
Plotting results (e.g., multipoles, angular correlations)

\item {} 
\sphinxAtStartPar
Comparing results across different cosmologies

\item {} 
\sphinxAtStartPar
Converting outputs into formats suitable for further analysis

\end{itemize}

\sphinxAtStartPar
\sphinxstylestrong{Notes}
\begin{itemize}
\item {} 
\sphinxAtStartPar
Pre\sphinxhyphen{}processing is external to \sphinxstylestrong{wlcf}; the code assumes inputs are already prepared.

\item {} 
\sphinxAtStartPar
Post\sphinxhyphen{}processing is flexible and depends on the scientific goals.

\item {} 
\sphinxAtStartPar
For large parameter scans, combining the Python interface with HPC resources is recommended.

\end{itemize}

\sphinxstepscope


\chapter{AddOn’s}
\label{\detokenize{addons:addon-s}}\label{\detokenize{addons::doc}}
\sphinxAtStartPar
This section describes how to extend the functionality of \sphinxstylestrong{wlcf} by adding new features, models, or utilities to the codebase.

\sphinxAtStartPar
The modular structure of the code allows users to include additional functionality through custom header and source files.

\sphinxAtStartPar
\sphinxstylestrong{Include files}

\sphinxAtStartPar
Header files define the interface of new functionalities and must be placed in the appropriate include directories.

\sphinxAtStartPar
To add new features:
\begin{enumerate}
\sphinxsetlistlabels{\arabic}{enumi}{enumii}{}{.}%
\item {} 
\sphinxAtStartPar
Create a header file (e.g., \sphinxtitleref{my\_module.h})

\item {} 
\sphinxAtStartPar
Place it in one of the include directories:

\begin{sphinxVerbatim}[commandchars=\\\{\}]
\PYG{n}{include}\PYG{o}{/}
\PYG{n}{addons}\PYG{o}{/}\PYG{n}{addons\PYGZus{}include}\PYG{o}{/}
\end{sphinxVerbatim}

\item {} 
\sphinxAtStartPar
Declare functions, structures, or parameters inside the header.

\end{enumerate}

\sphinxAtStartPar
Example

\sphinxAtStartPar
c:

\begin{sphinxVerbatim}[commandchars=\\\{\}]
\PYG{c+c1}{\PYGZsh{}ifndef MY\PYGZus{}MODULE\PYGZus{}H}
\PYG{c+c1}{\PYGZsh{}define MY\PYGZus{}MODULE\PYGZus{}H}
\PYG{n}{void} \PYG{n}{my\PYGZus{}new\PYGZus{}function}\PYG{p}{(}\PYG{n}{double} \PYG{o}{*}\PYG{n+nb}{input}\PYG{p}{,} \PYG{n}{double} \PYG{o}{*}\PYG{n}{output}\PYG{p}{)}\PYG{p}{;}
\PYG{c+c1}{\PYGZsh{}endif}
\end{sphinxVerbatim}
\begin{enumerate}
\sphinxsetlistlabels{\arabic}{enumi}{enumii}{}{.}%
\setcounter{enumi}{3}
\item {} 
\sphinxAtStartPar
Include the header in the relevant source files:

\end{enumerate}

\sphinxAtStartPar
c:

\begin{sphinxVerbatim}[commandchars=\\\{\}]
\PYG{c+c1}{\PYGZsh{}include \PYGZdq{}my\PYGZus{}module.h\PYGZdq{}}
\end{sphinxVerbatim}

\sphinxAtStartPar
\sphinxstylestrong{Source files}

\sphinxAtStartPar
Source files implement the functionality declared in the header files.

\sphinxAtStartPar
To add a new module:
\begin{enumerate}
\sphinxsetlistlabels{\arabic}{enumi}{enumii}{}{.}%
\item {} 
\sphinxAtStartPar
Create a source file (e.g., \sphinxtitleref{my\_module.c})

\item {} 
\sphinxAtStartPar
Place it in one of the source directories:

\begin{sphinxVerbatim}[commandchars=\\\{\}]
\PYG{n}{source}\PYG{o}{/}
\PYG{n}{addons}\PYG{o}{/}
\end{sphinxVerbatim}

\item {} 
\sphinxAtStartPar
Implement the functions:

\end{enumerate}

\sphinxAtStartPar
c:

\begin{sphinxVerbatim}[commandchars=\\\{\}]
\PYG{c+c1}{\PYGZsh{}include \PYGZdq{}my\PYGZus{}module.h\PYGZdq{}}
\PYG{n}{void} \PYG{n}{my\PYGZus{}new\PYGZus{}function}\PYG{p}{(}\PYG{n}{double} \PYG{o}{*}\PYG{n+nb}{input}\PYG{p}{,} \PYG{n}{double} \PYG{o}{*}\PYG{n}{output}\PYG{p}{)} \PYG{p}{\PYGZob{}}
\PYG{o}{/}\PYG{o}{/} \PYG{n}{Example} \PYG{n}{computation}
\PYG{o}{*}\PYG{n}{output} \PYG{o}{=} \PYG{n+nb}{input}\PYG{p}{[}\PYG{l+m+mi}{0}\PYG{p}{]} \PYG{o}{*} \PYG{l+m+mf}{2.0}\PYG{p}{;}\PYG{p}{\PYGZcb{}}
\end{sphinxVerbatim}

\sphinxAtStartPar
\sphinxstylestrong{Compilation}

\sphinxAtStartPar
To ensure the new files are compiled:
\begin{itemize}
\item {} 
\sphinxAtStartPar
Add the source file to the build system (if not automatically included)

\item {} 
\sphinxAtStartPar
Modify \sphinxtitleref{Makefile} or \sphinxtitleref{Makefile\_machine} if necessary

\end{itemize}

\sphinxAtStartPar
Example (if manual inclusion is needed):

\sphinxAtStartPar
make:

\begin{sphinxVerbatim}[commandchars=\\\{\}]
\PYG{n}{SRC} \PYG{o}{+}\PYG{o}{=} \PYG{n}{source}\PYG{o}{/}\PYG{n}{my\PYGZus{}module}\PYG{o}{.}\PYG{n}{c}
\end{sphinxVerbatim}

\sphinxAtStartPar
\sphinxstylestrong{Integration}

\sphinxAtStartPar
To integrate the new functionality into \sphinxstylestrong{wlcf}:
\begin{itemize}
\item {} 
\sphinxAtStartPar
Call your new functions from existing modules (e.g., in \sphinxtitleref{procedures.c})

\item {} 
\sphinxAtStartPar
Ensure parameters are passed correctly

\item {} 
\sphinxAtStartPar
Update relevant headers if new parameters are introduced

\end{itemize}

\sphinxstepscope


\chapter{Python interface}
\label{\detokenize{python:python-interface}}\label{\detokenize{python::doc}}
\sphinxAtStartPar
This section describes how to interact with \sphinxstylestrong{wlcf} using Python, either through the command line interface or via the Cython wrapper.

\sphinxAtStartPar
\sphinxstylestrong{Using command line interface}

\sphinxAtStartPar
The simplest way to run \sphinxstylestrong{wlcf} from Python is by calling the executable through the system shell.

\sphinxAtStartPar
Example

\sphinxAtStartPar
python:

\begin{sphinxVerbatim}[commandchars=\\\{\}]
\PYG{k+kn}{import}\PYG{+w}{ }\PYG{n+nn}{os}

\PYG{n}{os}\PYG{o}{.}\PYG{n}{system}\PYG{p}{(}\PYG{l+s+s2}{\PYGZdq{}}\PYG{l+s+s2}{../wlcf z=0.5 sigma8=0.8 fnamePS=./input/linear\PYGZus{}pk.txt}\PYG{l+s+s2}{\PYGZdq{}}\PYG{p}{)}
\end{sphinxVerbatim}

\sphinxAtStartPar
Alternatively, using \sphinxtitleref{subprocess} :

\sphinxAtStartPar
python:

\begin{sphinxVerbatim}[commandchars=\\\{\}]
\PYG{k+kn}{import}\PYG{+w}{ }\PYG{n+nn}{subprocess}

\PYG{n}{subprocess}\PYG{o}{.}\PYG{n}{run}\PYG{p}{(}\PYG{p}{[}
    \PYG{l+s+s2}{\PYGZdq{}}\PYG{l+s+s2}{../wlcf}\PYG{l+s+s2}{\PYGZdq{}}\PYG{p}{,}
    \PYG{l+s+s2}{\PYGZdq{}}\PYG{l+s+s2}{z=0.5}\PYG{l+s+s2}{\PYGZdq{}}\PYG{p}{,}
    \PYG{l+s+s2}{\PYGZdq{}}\PYG{l+s+s2}{sigma8=0.8}\PYG{l+s+s2}{\PYGZdq{}}\PYG{p}{,}
    \PYG{l+s+s2}{\PYGZdq{}}\PYG{l+s+s2}{fnamePS=./input/linear\PYGZus{}pk.txt}\PYG{l+s+s2}{\PYGZdq{}}\PYG{p}{]}\PYG{p}{)}
\end{sphinxVerbatim}

\sphinxAtStartPar
This approach allows you to:
\begin{itemize}
\item {} 
\sphinxAtStartPar
run parameter scans

\item {} 
\sphinxAtStartPar
automate workflows

\item {} 
\sphinxAtStartPar
integrate wlcf into larger pipelines

\end{itemize}

\sphinxAtStartPar
When using the command line interface, parameters must be passed as \sphinxtitleref{key=value} without spaces.

\sphinxAtStartPar
\sphinxstylestrong{Cython interface}

\sphinxAtStartPar
The code provides a \sphinxstylestrong{Cython wrapper} that allows direct interaction with wlcf from Python without invoking the executable.

\sphinxAtStartPar
example

\sphinxAtStartPar
python:

\begin{sphinxVerbatim}[commandchars=\\\{\}]
\PYG{k+kn}{from}\PYG{+w}{ }\PYG{n+nn}{wlcfpy}\PYG{+w}{ }\PYG{k+kn}{import} \PYG{n}{wlcf}

\PYG{n}{w} \PYG{o}{=} \PYG{n}{wlcf}\PYG{p}{(}\PYG{p}{)}

\PYG{n}{w}\PYG{o}{.}\PYG{n}{set}\PYG{p}{(}
    \PYG{n}{numberThreads}\PYG{o}{=}\PYG{l+m+mi}{4}\PYG{p}{,}
    \PYG{n}{verbose}\PYG{o}{=}\PYG{l+m+mi}{1}\PYG{p}{,}
    \PYG{n}{fnamePS}\PYG{o}{=}\PYG{l+s+s2}{\PYGZdq{}}\PYG{l+s+s2}{./input/linear\PYGZus{}pk.txt}\PYG{l+s+s2}{\PYGZdq{}}\PYG{p}{)}

\PYG{n}{w}\PYG{o}{.}\PYG{n}{Run}\PYG{p}{(}\PYG{p}{)}
\end{sphinxVerbatim}

\sphinxAtStartPar
\sphinxstylestrong{Notes}
\begin{itemize}
\item {} 
\sphinxAtStartPar
The working directory should contain the required input files (e.g., power spectrum).

\item {} 
\sphinxAtStartPar
Output files are written in the same way as the C executable.

\item {} 
\sphinxAtStartPar
Some parameters (e.g., number of threads) may be controlled both via Python and internally by the code.

\end{itemize}

\sphinxstepscope


\chapter{Hands on}
\label{\detokenize{handson:hands-on}}\label{\detokenize{handson::doc}}
\sphinxAtStartPar
This section provides a practical example of how to use \sphinxstylestrong{wlcf} from a
Python notebook. The example follows a complete workflow:
\begin{enumerate}
\sphinxsetlistlabels{\arabic}{enumi}{enumii}{}{.}%
\item {} 
\sphinxAtStartPar
Generate a linear matter power spectrum using CAMB.

\item {} 
\sphinxAtStartPar
Run \sphinxcode{\sphinxupquote{wlcf}} through the Python wrapper.

\item {} 
\sphinxAtStartPar
Plot the resulting 3\sphinxhyphen{}point correlation function multipoles.

\item {} 
\sphinxAtStartPar
Compare different theoretical models.

\end{enumerate}

\sphinxAtStartPar
The notebook associated with this tutorial can be used as a starting point
for interactive runs and parameter exploration.


\section{Setting up the working directory}
\label{\detokenize{handson:setting-up-the-working-directory}}
\sphinxAtStartPar
First, import the required Python packages and move to the \sphinxcode{\sphinxupquote{tests}} directory
of the repository:

\begin{sphinxVerbatim}[commandchars=\\\{\}]
\PYG{k+kn}{import}\PYG{+w}{ }\PYG{n+nn}{os}
\PYG{k+kn}{import}\PYG{+w}{ }\PYG{n+nn}{numpy}\PYG{+w}{ }\PYG{k}{as}\PYG{+w}{ }\PYG{n+nn}{np}
\PYG{k+kn}{import}\PYG{+w}{ }\PYG{n+nn}{matplotlib}\PYG{n+nn}{.}\PYG{n+nn}{pyplot}\PYG{+w}{ }\PYG{k}{as}\PYG{+w}{ }\PYG{n+nn}{plt}

\PYG{k+kn}{import}\PYG{+w}{ }\PYG{n+nn}{camb}
\PYG{k+kn}{from}\PYG{+w}{ }\PYG{n+nn}{camb}\PYG{+w}{ }\PYG{k+kn}{import} \PYG{n}{model}
\PYG{k+kn}{from}\PYG{+w}{ }\PYG{n+nn}{wlcfpy}\PYG{+w}{ }\PYG{k+kn}{import} \PYG{n}{wlcf}

\PYG{n}{WLCF\PYGZus{}TESTS\PYGZus{}DIR} \PYG{o}{=} \PYG{l+s+s2}{\PYGZdq{}}\PYG{l+s+s2}{/path/to/wlcf/tests}\PYG{l+s+s2}{\PYGZdq{}}
\PYG{n}{os}\PYG{o}{.}\PYG{n}{chdir}\PYG{p}{(}\PYG{n}{WLCF\PYGZus{}TESTS\PYGZus{}DIR}\PYG{p}{)}

\PYG{n}{os}\PYG{o}{.}\PYG{n}{makedirs}\PYG{p}{(}\PYG{l+s+s2}{\PYGZdq{}}\PYG{l+s+s2}{input}\PYG{l+s+s2}{\PYGZdq{}}\PYG{p}{,} \PYG{n}{exist\PYGZus{}ok}\PYG{o}{=}\PYG{k+kc}{True}\PYG{p}{)}
\PYG{n}{os}\PYG{o}{.}\PYG{n}{makedirs}\PYG{p}{(}\PYG{l+s+s2}{\PYGZdq{}}\PYG{l+s+s2}{Output}\PYG{l+s+s2}{\PYGZdq{}}\PYG{p}{,} \PYG{n}{exist\PYGZus{}ok}\PYG{o}{=}\PYG{k+kc}{True}\PYG{p}{)}
\PYG{n}{os}\PYG{o}{.}\PYG{n}{makedirs}\PYG{p}{(}\PYG{l+s+s2}{\PYGZdq{}}\PYG{l+s+s2}{Bell\PYGZus{}outputs}\PYG{l+s+s2}{\PYGZdq{}}\PYG{p}{,} \PYG{n}{exist\PYGZus{}ok}\PYG{o}{=}\PYG{k+kc}{True}\PYG{p}{)}
\end{sphinxVerbatim}


\section{Generating the linear power spectrum}
\label{\detokenize{handson:generating-the-linear-power-spectrum}}
\sphinxAtStartPar
The linear matter power spectrum can be generated with CAMB. In this example
we use a T17\sphinxhyphen{}like cosmology at redshift \(z = 0.5078\):

\begin{sphinxVerbatim}[commandchars=\\\{\}]
\PYG{k}{def}\PYG{+w}{ }\PYG{n+nf}{make\PYGZus{}camb\PYGZus{}linear\PYGZus{}pk}\PYG{p}{(}
    \PYG{n}{output\PYGZus{}file}\PYG{o}{=}\PYG{l+s+s2}{\PYGZdq{}}\PYG{l+s+s2}{./input/linear\PYGZus{}pk\PYGZus{}camb\PYGZus{}T17\PYGZus{}z05078.txt}\PYG{l+s+s2}{\PYGZdq{}}\PYG{p}{,}
    \PYG{n}{z\PYGZus{}pk}\PYG{o}{=}\PYG{l+m+mf}{0.5078}\PYG{p}{,}
    \PYG{n}{h}\PYG{o}{=}\PYG{l+m+mf}{0.7}\PYG{p}{,}
    \PYG{n}{Omb}\PYG{o}{=}\PYG{l+m+mf}{0.046}\PYG{p}{,}
    \PYG{n}{Omc}\PYG{o}{=}\PYG{l+m+mf}{0.233}\PYG{p}{,}
    \PYG{n}{Omnu}\PYG{o}{=}\PYG{l+m+mf}{0.0}\PYG{p}{,}
    \PYG{n}{ns}\PYG{o}{=}\PYG{l+m+mf}{0.97}\PYG{p}{,}
    \PYG{n}{sigma8}\PYG{o}{=}\PYG{l+m+mf}{0.82}\PYG{p}{,}
    \PYG{n}{w0}\PYG{o}{=}\PYG{o}{\PYGZhy{}}\PYG{l+m+mf}{1.0}\PYG{p}{,}
    \PYG{n}{kmin}\PYG{o}{=}\PYG{l+m+mf}{1e\PYGZhy{}4}\PYG{p}{,}
    \PYG{n}{kmax}\PYG{o}{=}\PYG{l+m+mf}{50.0}\PYG{p}{,}
    \PYG{n}{npoints}\PYG{o}{=}\PYG{l+m+mi}{1000}\PYG{p}{,}
\PYG{p}{)}\PYG{p}{:}

    \PYG{n}{os}\PYG{o}{.}\PYG{n}{makedirs}\PYG{p}{(}\PYG{n}{os}\PYG{o}{.}\PYG{n}{path}\PYG{o}{.}\PYG{n}{dirname}\PYG{p}{(}\PYG{n}{output\PYGZus{}file}\PYG{p}{)}\PYG{p}{,} \PYG{n}{exist\PYGZus{}ok}\PYG{o}{=}\PYG{k+kc}{True}\PYG{p}{)}

    \PYG{n}{H0} \PYG{o}{=} \PYG{l+m+mf}{100.0} \PYG{o}{*} \PYG{n}{h}
    \PYG{n}{ombh2} \PYG{o}{=} \PYG{n}{Omb} \PYG{o}{*} \PYG{n}{h}\PYG{o}{*}\PYG{o}{*}\PYG{l+m+mi}{2}
    \PYG{n}{omch2} \PYG{o}{=} \PYG{n}{Omc} \PYG{o}{*} \PYG{n}{h}\PYG{o}{*}\PYG{o}{*}\PYG{l+m+mi}{2}

    \PYG{n}{pars} \PYG{o}{=} \PYG{n}{camb}\PYG{o}{.}\PYG{n}{CAMBparams}\PYG{p}{(}\PYG{p}{)}

    \PYG{c+c1}{\PYGZsh{} For now we assume massless neutrinos, consistent with Omnu = 0.}
    \PYG{n}{pars}\PYG{o}{.}\PYG{n}{set\PYGZus{}cosmology}\PYG{p}{(}
        \PYG{n}{H0}\PYG{o}{=}\PYG{n}{H0}\PYG{p}{,}
        \PYG{n}{ombh2}\PYG{o}{=}\PYG{n}{ombh2}\PYG{p}{,}
        \PYG{n}{omch2}\PYG{o}{=}\PYG{n}{omch2}\PYG{p}{,}
    \PYG{p}{)}

    \PYG{n}{pars}\PYG{o}{.}\PYG{n}{InitPower}\PYG{o}{.}\PYG{n}{set\PYGZus{}params}\PYG{p}{(}
        \PYG{n}{ns}\PYG{o}{=}\PYG{n}{ns}\PYG{p}{,}
        \PYG{n}{As}\PYG{o}{=}\PYG{l+m+mf}{2.1e\PYGZhy{}9}
    \PYG{p}{)}

    \PYG{n}{pars}\PYG{o}{.}\PYG{n}{set\PYGZus{}dark\PYGZus{}energy}\PYG{p}{(}\PYG{n}{w}\PYG{o}{=}\PYG{n}{w0}\PYG{p}{)}

    \PYG{c+c1}{\PYGZsh{} Linear power spectrum at the same redshift used by wlcf/zbin}
    \PYG{n}{pars}\PYG{o}{.}\PYG{n}{set\PYGZus{}matter\PYGZus{}power}\PYG{p}{(}
        \PYG{n}{redshifts}\PYG{o}{=}\PYG{p}{[}\PYG{n}{z\PYGZus{}pk}\PYG{p}{]}\PYG{p}{,}
        \PYG{n}{kmax}\PYG{o}{=}\PYG{n}{kmax}
    \PYG{p}{)}

    \PYG{n}{pars}\PYG{o}{.}\PYG{n}{NonLinear} \PYG{o}{=} \PYG{n}{model}\PYG{o}{.}\PYG{n}{NonLinear\PYGZus{}none}

    \PYG{c+c1}{\PYGZsh{} First run to compute sigma8 for the trial As}
    \PYG{n}{results} \PYG{o}{=} \PYG{n}{camb}\PYG{o}{.}\PYG{n}{get\PYGZus{}results}\PYG{p}{(}\PYG{n}{pars}\PYG{p}{)}
    \PYG{n}{sigma8\PYGZus{}current} \PYG{o}{=} \PYG{n}{results}\PYG{o}{.}\PYG{n}{get\PYGZus{}sigma8}\PYG{p}{(}\PYG{p}{)}\PYG{p}{[}\PYG{l+m+mi}{0}\PYG{p}{]}

    \PYG{c+c1}{\PYGZsh{} Rescale As so that CAMB matches the requested sigma8}
    \PYG{n}{As\PYGZus{}rescaled} \PYG{o}{=} \PYG{l+m+mf}{2.1e\PYGZhy{}9} \PYG{o}{*} \PYG{p}{(}\PYG{n}{sigma8} \PYG{o}{/} \PYG{n}{sigma8\PYGZus{}current}\PYG{p}{)} \PYG{o}{*}\PYG{o}{*} \PYG{l+m+mi}{2}

    \PYG{n}{pars}\PYG{o}{.}\PYG{n}{InitPower}\PYG{o}{.}\PYG{n}{set\PYGZus{}params}\PYG{p}{(}
        \PYG{n}{ns}\PYG{o}{=}\PYG{n}{ns}\PYG{p}{,}
        \PYG{n}{As}\PYG{o}{=}\PYG{n}{As\PYGZus{}rescaled}
    \PYG{p}{)}

    \PYG{n}{results} \PYG{o}{=} \PYG{n}{camb}\PYG{o}{.}\PYG{n}{get\PYGZus{}results}\PYG{p}{(}\PYG{n}{pars}\PYG{p}{)}

    \PYG{n}{kh}\PYG{p}{,} \PYG{n}{redshifts}\PYG{p}{,} \PYG{n}{pk} \PYG{o}{=} \PYG{n}{results}\PYG{o}{.}\PYG{n}{get\PYGZus{}matter\PYGZus{}power\PYGZus{}spectrum}\PYG{p}{(}
        \PYG{n}{minkh}\PYG{o}{=}\PYG{n}{kmin}\PYG{p}{,}
        \PYG{n}{maxkh}\PYG{o}{=}\PYG{n}{kmax}\PYG{p}{,}
        \PYG{n}{npoints}\PYG{o}{=}\PYG{n}{npoints}\PYG{p}{,}
    \PYG{p}{)}

    \PYG{n}{pk\PYGZus{}z} \PYG{o}{=} \PYG{n}{pk}\PYG{p}{[}\PYG{l+m+mi}{0}\PYG{p}{]}

    \PYG{n}{np}\PYG{o}{.}\PYG{n}{savetxt}\PYG{p}{(}
        \PYG{n}{output\PYGZus{}file}\PYG{p}{,}
        \PYG{n}{np}\PYG{o}{.}\PYG{n}{column\PYGZus{}stack}\PYG{p}{(}\PYG{p}{[}\PYG{n}{kh}\PYG{p}{,} \PYG{n}{pk\PYGZus{}z}\PYG{p}{]}\PYG{p}{)}\PYG{p}{,}
        \PYG{n}{fmt}\PYG{o}{=}\PYG{l+s+s2}{\PYGZdq{}}\PYG{l+s+si}{\PYGZpc{}.10e}\PYG{l+s+s2}{\PYGZdq{}}
    \PYG{p}{)}

    \PYG{n+nb}{print}\PYG{p}{(}\PYG{l+s+sa}{f}\PYG{l+s+s2}{\PYGZdq{}}\PYG{l+s+s2}{Saved linear P(k) at z=}\PYG{l+s+si}{\PYGZob{}}\PYG{n}{z\PYGZus{}pk}\PYG{l+s+si}{\PYGZcb{}}\PYG{l+s+s2}{ to:}\PYG{l+s+s2}{\PYGZdq{}}\PYG{p}{)}
    \PYG{n+nb}{print}\PYG{p}{(}\PYG{n}{output\PYGZus{}file}\PYG{p}{)}
    \PYG{n+nb}{print}\PYG{p}{(}\PYG{l+s+sa}{f}\PYG{l+s+s2}{\PYGZdq{}}\PYG{l+s+s2}{sigma8 target  = }\PYG{l+s+si}{\PYGZob{}}\PYG{n}{sigma8}\PYG{l+s+si}{\PYGZcb{}}\PYG{l+s+s2}{\PYGZdq{}}\PYG{p}{)}
    \PYG{n+nb}{print}\PYG{p}{(}\PYG{l+s+sa}{f}\PYG{l+s+s2}{\PYGZdq{}}\PYG{l+s+s2}{sigma8 CAMB    = }\PYG{l+s+si}{\PYGZob{}}\PYG{n}{results}\PYG{o}{.}\PYG{n}{get\PYGZus{}sigma8}\PYG{p}{(}\PYG{p}{)}\PYG{p}{[}\PYG{l+m+mi}{0}\PYG{p}{]}\PYG{l+s+si}{:}\PYG{l+s+s2}{.6f}\PYG{l+s+si}{\PYGZcb{}}\PYG{l+s+s2}{\PYGZdq{}}\PYG{p}{)}
    \PYG{n+nb}{print}\PYG{p}{(}\PYG{l+s+sa}{f}\PYG{l+s+s2}{\PYGZdq{}}\PYG{l+s+s2}{As rescaled    = }\PYG{l+s+si}{\PYGZob{}}\PYG{n}{As\PYGZus{}rescaled}\PYG{l+s+si}{:}\PYG{l+s+s2}{.6e}\PYG{l+s+si}{\PYGZcb{}}\PYG{l+s+s2}{\PYGZdq{}}\PYG{p}{)}

    \PYG{k}{return} \PYG{n}{output\PYGZus{}file}\PYG{p}{,} \PYG{n}{kh}\PYG{p}{,} \PYG{n}{pk\PYGZus{}z}
\end{sphinxVerbatim}

\sphinxAtStartPar
The resulting file is passed to \sphinxcode{\sphinxupquote{wlcf}} through the parameter \sphinxcode{\sphinxupquote{fnamePS}}.


\section{Running wlcf from Python}
\label{\detokenize{handson:running-wlcf-from-python}}
\sphinxAtStartPar
The code can be executed directly from Python using the \sphinxcode{\sphinxupquote{wlcfpy}} wrapper:

\begin{sphinxVerbatim}[commandchars=\\\{\}]
\PYG{n}{w} \PYG{o}{=} \PYG{n}{wlcf}\PYG{p}{(}\PYG{p}{)}
\PYG{n}{w}\PYG{o}{.}\PYG{n}{set}\PYG{p}{(}
    \PYG{n}{numberThreads}\PYG{o}{=}\PYG{l+m+mi}{8}\PYG{p}{,}
    \PYG{n}{verbose}\PYG{o}{=}\PYG{l+m+mi}{2}\PYG{p}{,}
    \PYG{n}{verbose\PYGZus{}log}\PYG{o}{=}\PYG{l+m+mi}{1}\PYG{p}{,}

    \PYG{n}{z}\PYG{o}{=}\PYG{l+m+mf}{0.5078}\PYG{p}{,}
    \PYG{n}{h}\PYG{o}{=}\PYG{l+m+mf}{0.7}\PYG{p}{,}
    \PYG{n}{sigma8}\PYG{o}{=}\PYG{l+m+mf}{0.82}\PYG{p}{,}
    \PYG{n}{Omb}\PYG{o}{=}\PYG{l+m+mf}{0.046}\PYG{p}{,}
    \PYG{n}{Omc}\PYG{o}{=}\PYG{l+m+mf}{0.233}\PYG{p}{,}
    \PYG{n}{Omnu}\PYG{o}{=}\PYG{l+m+mf}{0.0}\PYG{p}{,}
    \PYG{n}{ns}\PYG{o}{=}\PYG{l+m+mf}{0.97}\PYG{p}{,}
    \PYG{n}{w}\PYG{o}{=}\PYG{o}{\PYGZhy{}}\PYG{l+m+mf}{1.0}\PYG{p}{,}

    \PYG{n}{fnamePS}\PYG{o}{=}\PYG{l+s+s2}{\PYGZdq{}}\PYG{l+s+s2}{./input/linear\PYGZus{}pk\PYGZus{}camb\PYGZus{}T17\PYGZus{}z05078.txt}\PYG{l+s+s2}{\PYGZdq{}}\PYG{p}{,}
    \PYG{n}{Wg}\PYG{o}{=}\PYG{l+m+mi}{0}\PYG{p}{,}
    \PYG{n}{fWgchi}\PYG{o}{=}\PYG{l+s+s2}{\PYGZdq{}}\PYG{l+s+s2}{./input/Wg\PYGZus{}Takahashi\PYGZus{}z05078.txt}\PYG{l+s+s2}{\PYGZdq{}}\PYG{p}{,}

    \PYG{n}{rootDir}\PYG{o}{=}\PYG{l+s+s2}{\PYGZdq{}}\PYG{l+s+s2}{Output}\PYG{l+s+s2}{\PYGZdq{}}\PYG{p}{,}
    \PYG{n}{path\PYGZus{}Bells}\PYG{o}{=}\PYG{l+s+s2}{\PYGZdq{}}\PYG{l+s+s2}{Bell\PYGZus{}outputs}\PYG{l+s+s2}{\PYGZdq{}}\PYG{p}{,}
    \PYG{n}{prefix}\PYG{o}{=}\PYG{l+s+s2}{\PYGZdq{}}\PYG{l+s+s2}{halo\PYGZus{}model\PYGZus{}z05078\PYGZus{}}\PYG{l+s+s2}{\PYGZdq{}}\PYG{p}{,}

    \PYG{n}{tree\PYGZus{}level}\PYG{o}{=}\PYG{l+m+mi}{4}\PYG{p}{,}
    \PYG{n}{zbin}\PYG{o}{=}\PYG{l+m+mf}{0.5078}\PYG{p}{,}
    \PYG{n}{mMax}\PYG{o}{=}\PYG{l+m+mi}{5}\PYG{p}{,}
    \PYG{n}{chiQuadSteps}\PYG{o}{=}\PYG{l+m+mi}{300}\PYG{p}{,}
    \PYG{n}{GLpoints}\PYG{o}{=}\PYG{l+m+mi}{64}\PYG{p}{,}
    \PYG{n}{Nell}\PYG{o}{=}\PYG{l+m+mi}{128}\PYG{p}{,}
    \PYG{n}{ellmax}\PYG{o}{=}\PYG{l+m+mi}{10000}\PYG{p}{,}
    \PYG{n}{ellmin}\PYG{o}{=}\PYG{l+m+mf}{0.001}\PYG{p}{,}
    \PYG{n}{writevectors}\PYG{o}{=}\PYG{l+m+mi}{0}\PYG{p}{,}
    \PYG{n}{chatty}\PYG{o}{=}\PYG{l+m+mi}{2}\PYG{p}{,}
    \PYG{n}{options}\PYG{o}{=}\PYG{l+s+s2}{\PYGZdq{}}\PYG{l+s+s2}{\PYGZdq{}}
\PYG{p}{)}

\PYG{n}{w}\PYG{o}{.}\PYG{n}{Run}\PYG{p}{(}\PYG{p}{)}
\end{sphinxVerbatim}

\sphinxAtStartPar
Here \sphinxcode{\sphinxupquote{tree\_level=4}} selects the Takahashi/Halo\sphinxhyphen{}model branch.


\section{Plotting the 3PCF multipoles}
\label{\detokenize{handson:plotting-the-3pcf-multipoles}}
\sphinxAtStartPar
The output files \sphinxcode{\sphinxupquote{zetam*.txt}} contain the multipoles
\(\zeta_m(\theta_1,\theta_2)\). These can be visualized as 2D maps:

\begin{sphinxVerbatim}[commandchars=\\\{\}]
\PYG{k+kn}{import}\PYG{+w}{ }\PYG{n+nn}{numpy}\PYG{+w}{ }\PYG{k}{as}\PYG{+w}{ }\PYG{n+nn}{np}
\PYG{k+kn}{import}\PYG{+w}{ }\PYG{n+nn}{matplotlib}\PYG{n+nn}{.}\PYG{n+nn}{pyplot}\PYG{+w}{ }\PYG{k}{as}\PYG{+w}{ }\PYG{n+nn}{plt}
\PYG{k+kn}{from}\PYG{+w}{ }\PYG{n+nn}{matplotlib}\PYG{n+nn}{.}\PYG{n+nn}{colors}\PYG{+w}{ }\PYG{k+kn}{import} \PYG{n}{LogNorm}

\PYG{n}{path} \PYG{o}{=} \PYG{l+s+s2}{\PYGZdq{}}\PYG{l+s+s2}{/home/sadi/LSST/test/wlcf\PYGZhy{}main/tests/Bell\PYGZus{}outputs}\PYG{l+s+s2}{\PYGZdq{}}
\PYG{n}{prefix} \PYG{o}{=} \PYG{l+s+s2}{\PYGZdq{}}\PYG{l+s+s2}{halo\PYGZus{}model\PYGZus{}z05078\PYGZus{}}\PYG{l+s+s2}{\PYGZdq{}}

\PYG{n}{theta} \PYG{o}{=} \PYG{n}{np}\PYG{o}{.}\PYG{n}{loadtxt}\PYG{p}{(}\PYG{l+s+sa}{f}\PYG{l+s+s2}{\PYGZdq{}}\PYG{l+s+si}{\PYGZob{}}\PYG{n}{path}\PYG{l+s+si}{\PYGZcb{}}\PYG{l+s+s2}{/}\PYG{l+s+si}{\PYGZob{}}\PYG{n}{prefix}\PYG{l+s+si}{\PYGZcb{}}\PYG{l+s+s2}{theta\PYGZus{}array.txt}\PYG{l+s+s2}{\PYGZdq{}}\PYG{p}{)} \PYG{o}{*} \PYG{l+m+mi}{180} \PYG{o}{/} \PYG{n}{np}\PYG{o}{.}\PYG{n}{pi} \PYG{o}{*} \PYG{l+m+mi}{60}
\PYG{n}{Theta2}\PYG{p}{,} \PYG{n}{Theta1} \PYG{o}{=} \PYG{n}{np}\PYG{o}{.}\PYG{n}{meshgrid}\PYG{p}{(}\PYG{n}{theta}\PYG{p}{,} \PYG{n}{theta}\PYG{p}{)}

\PYG{n}{fig} \PYG{o}{=} \PYG{n}{plt}\PYG{o}{.}\PYG{n}{figure}\PYG{p}{(}\PYG{n}{figsize}\PYG{o}{=}\PYG{p}{(}\PYG{l+m+mi}{19}\PYG{p}{,} \PYG{l+m+mi}{4}\PYG{p}{)}\PYG{p}{)}

\PYG{n}{gs} \PYG{o}{=} \PYG{n}{fig}\PYG{o}{.}\PYG{n}{add\PYGZus{}gridspec}\PYG{p}{(}
    \PYG{n}{nrows}\PYG{o}{=}\PYG{l+m+mi}{1}\PYG{p}{,}
    \PYG{n}{ncols}\PYG{o}{=}\PYG{l+m+mi}{6}\PYG{p}{,}
    \PYG{n}{width\PYGZus{}ratios}\PYG{o}{=}\PYG{p}{[}\PYG{l+m+mi}{1}\PYG{p}{,} \PYG{l+m+mi}{1}\PYG{p}{,} \PYG{l+m+mi}{1}\PYG{p}{,} \PYG{l+m+mi}{1}\PYG{p}{,} \PYG{l+m+mi}{1}\PYG{p}{,} \PYG{l+m+mf}{0.06}\PYG{p}{]}\PYG{p}{,}
    \PYG{n}{wspace}\PYG{o}{=}\PYG{l+m+mf}{0.25}
\PYG{p}{)}

\PYG{n}{norm} \PYG{o}{=} \PYG{n}{LogNorm}\PYG{p}{(}\PYG{n}{vmin}\PYG{o}{=}\PYG{l+m+mf}{1e\PYGZhy{}11}\PYG{p}{,} \PYG{n}{vmax}\PYG{o}{=}\PYG{l+m+mf}{1e\PYGZhy{}8}\PYG{p}{)}

\PYG{k}{for} \PYG{n}{m} \PYG{o+ow}{in} \PYG{n+nb}{range}\PYG{p}{(}\PYG{l+m+mi}{5}\PYG{p}{)}\PYG{p}{:}
    \PYG{n}{ax} \PYG{o}{=} \PYG{n}{fig}\PYG{o}{.}\PYG{n}{add\PYGZus{}subplot}\PYG{p}{(}\PYG{n}{gs}\PYG{p}{[}\PYG{l+m+mi}{0}\PYG{p}{,} \PYG{n}{m}\PYG{p}{]}\PYG{p}{)}

    \PYG{n}{zeta} \PYG{o}{=} \PYG{n}{np}\PYG{o}{.}\PYG{n}{abs}\PYG{p}{(}\PYG{n}{np}\PYG{o}{.}\PYG{n}{loadtxt}\PYG{p}{(}\PYG{l+s+sa}{f}\PYG{l+s+s2}{\PYGZdq{}}\PYG{l+s+si}{\PYGZob{}}\PYG{n}{path}\PYG{l+s+si}{\PYGZcb{}}\PYG{l+s+s2}{/}\PYG{l+s+si}{\PYGZob{}}\PYG{n}{prefix}\PYG{l+s+si}{\PYGZcb{}}\PYG{l+s+s2}{zetam}\PYG{l+s+si}{\PYGZob{}}\PYG{n}{m}\PYG{l+s+si}{\PYGZcb{}}\PYG{l+s+s2}{.txt}\PYG{l+s+s2}{\PYGZdq{}}\PYG{p}{)}\PYG{p}{)}

    \PYG{n}{im} \PYG{o}{=} \PYG{n}{ax}\PYG{o}{.}\PYG{n}{pcolormesh}\PYG{p}{(}
        \PYG{n}{Theta2}\PYG{p}{,}
        \PYG{n}{Theta1}\PYG{p}{,}
        \PYG{n}{zeta}\PYG{p}{,}
        \PYG{n}{shading}\PYG{o}{=}\PYG{l+s+s2}{\PYGZdq{}}\PYG{l+s+s2}{auto}\PYG{l+s+s2}{\PYGZdq{}}\PYG{p}{,}
        \PYG{n}{cmap}\PYG{o}{=}\PYG{l+s+s2}{\PYGZdq{}}\PYG{l+s+s2}{RdYlBu\PYGZus{}r}\PYG{l+s+s2}{\PYGZdq{}}\PYG{p}{,}
        \PYG{n}{norm}\PYG{o}{=}\PYG{n}{norm}
    \PYG{p}{)}

    \PYG{n}{ax}\PYG{o}{.}\PYG{n}{set\PYGZus{}xscale}\PYG{p}{(}\PYG{l+s+s2}{\PYGZdq{}}\PYG{l+s+s2}{log}\PYG{l+s+s2}{\PYGZdq{}}\PYG{p}{)}
    \PYG{n}{ax}\PYG{o}{.}\PYG{n}{set\PYGZus{}yscale}\PYG{p}{(}\PYG{l+s+s2}{\PYGZdq{}}\PYG{l+s+s2}{log}\PYG{l+s+s2}{\PYGZdq{}}\PYG{p}{)}
    \PYG{n}{ax}\PYG{o}{.}\PYG{n}{set\PYGZus{}xlim}\PYG{p}{(}\PYG{l+m+mi}{10}\PYG{p}{,} \PYG{l+m+mi}{200}\PYG{p}{)}
    \PYG{n}{ax}\PYG{o}{.}\PYG{n}{set\PYGZus{}ylim}\PYG{p}{(}\PYG{l+m+mi}{10}\PYG{p}{,} \PYG{l+m+mi}{200}\PYG{p}{)}
    \PYG{n}{ax}\PYG{o}{.}\PYG{n}{set\PYGZus{}aspect}\PYG{p}{(}\PYG{l+s+s2}{\PYGZdq{}}\PYG{l+s+s2}{equal}\PYG{l+s+s2}{\PYGZdq{}}\PYG{p}{,} \PYG{n}{adjustable}\PYG{o}{=}\PYG{l+s+s2}{\PYGZdq{}}\PYG{l+s+s2}{box}\PYG{l+s+s2}{\PYGZdq{}}\PYG{p}{)}
    \PYG{n}{ax}\PYG{o}{.}\PYG{n}{set\PYGZus{}title}\PYG{p}{(}\PYG{l+s+sa}{fr}\PYG{l+s+s2}{\PYGZdq{}}\PYG{l+s+s2}{\PYGZdl{}m=}\PYG{l+s+si}{\PYGZob{}}\PYG{n}{m}\PYG{l+s+si}{\PYGZcb{}}\PYG{l+s+s2}{\PYGZdl{}}\PYG{l+s+s2}{\PYGZdq{}}\PYG{p}{,} \PYG{n}{fontsize}\PYG{o}{=}\PYG{l+m+mi}{14}\PYG{p}{)}

    \PYG{k}{if} \PYG{n}{m} \PYG{o}{==} \PYG{l+m+mi}{0}\PYG{p}{:}
        \PYG{n}{ax}\PYG{o}{.}\PYG{n}{set\PYGZus{}ylabel}\PYG{p}{(}\PYG{l+s+sa}{r}\PYG{l+s+s2}{\PYGZdq{}}\PYG{l+s+s2}{\PYGZdl{}}\PYG{l+s+s2}{\PYGZbs{}}\PYG{l+s+s2}{theta\PYGZus{}1\PYGZdl{} [arcmin]}\PYG{l+s+s2}{\PYGZdq{}}\PYG{p}{,} \PYG{n}{fontsize}\PYG{o}{=}\PYG{l+m+mi}{13}\PYG{p}{)}

    \PYG{n}{ax}\PYG{o}{.}\PYG{n}{set\PYGZus{}xlabel}\PYG{p}{(}\PYG{l+s+sa}{r}\PYG{l+s+s2}{\PYGZdq{}}\PYG{l+s+s2}{\PYGZdl{}}\PYG{l+s+s2}{\PYGZbs{}}\PYG{l+s+s2}{theta\PYGZus{}2\PYGZdl{} [arcmin]}\PYG{l+s+s2}{\PYGZdq{}}\PYG{p}{,} \PYG{n}{fontsize}\PYG{o}{=}\PYG{l+m+mi}{13}\PYG{p}{)}

\PYG{c+c1}{\PYGZsh{} Colorbar outside the plots}
\PYG{n}{cax} \PYG{o}{=} \PYG{n}{fig}\PYG{o}{.}\PYG{n}{add\PYGZus{}subplot}\PYG{p}{(}\PYG{n}{gs}\PYG{p}{[}\PYG{l+m+mi}{0}\PYG{p}{,} \PYG{o}{\PYGZhy{}}\PYG{l+m+mi}{1}\PYG{p}{]}\PYG{p}{)}
\PYG{n}{cbar} \PYG{o}{=} \PYG{n}{fig}\PYG{o}{.}\PYG{n}{colorbar}\PYG{p}{(}\PYG{n}{im}\PYG{p}{,} \PYG{n}{cax}\PYG{o}{=}\PYG{n}{cax}\PYG{p}{)}
\PYG{n}{cbar}\PYG{o}{.}\PYG{n}{set\PYGZus{}label}\PYG{p}{(}\PYG{l+s+sa}{r}\PYG{l+s+s2}{\PYGZdq{}}\PYG{l+s+s2}{\PYGZdl{}|}\PYG{l+s+s2}{\PYGZbs{}}\PYG{l+s+s2}{zeta\PYGZus{}m|\PYGZdl{}}\PYG{l+s+s2}{\PYGZdq{}}\PYG{p}{,} \PYG{n}{fontsize}\PYG{o}{=}\PYG{l+m+mi}{13}\PYG{p}{)}
\PYG{n}{cbar}\PYG{o}{.}\PYG{n}{ax}\PYG{o}{.}\PYG{n}{tick\PYGZus{}params}\PYG{p}{(}\PYG{n}{labelsize}\PYG{o}{=}\PYG{l+m+mi}{11}\PYG{p}{)}

\PYG{n}{fig}\PYG{o}{.}\PYG{n}{suptitle}\PYG{p}{(}
    \PYG{l+s+s2}{\PYGZdq{}}\PYG{l+s+s2}{Halo Model: 3PCF Multipoles at z = 0.5078}\PYG{l+s+s2}{\PYGZdq{}}\PYG{p}{,}
    \PYG{n}{fontsize}\PYG{o}{=}\PYG{l+m+mi}{18}\PYG{p}{,}
    \PYG{n}{y}\PYG{o}{=}\PYG{l+m+mf}{1.02}
\PYG{p}{)}
    \PYG{n}{plt}\PYG{o}{.}\PYG{n}{savefig}\PYG{p}{(}\PYG{l+s+s2}{\PYGZdq{}}\PYG{l+s+s2}{halo\PYGZus{}model\PYGZus{}3pcf.png}\PYG{l+s+s2}{\PYGZdq{}}\PYG{p}{,} \PYG{n}{dpi}\PYG{o}{=}\PYG{l+m+mi}{300}\PYG{p}{,} \PYG{n}{bbox\PYGZus{}inches}\PYG{o}{=}\PYG{l+s+s2}{\PYGZdq{}}\PYG{l+s+s2}{tight}\PYG{l+s+s2}{\PYGZdq{}}\PYG{p}{)}
    \PYG{n}{plt}\PYG{o}{.}\PYG{n}{show}\PYG{p}{(}\PYG{p}{)}
\end{sphinxVerbatim}

\noindent{\hspace*{\fill}\sphinxincludegraphics[width=1.000\linewidth]{{halo_model_3pcf}.png}\hspace*{\fill}}

\sphinxAtStartPar
Or we can plot

\begin{sphinxVerbatim}[commandchars=\\\{\}]
\PYG{n}{path} \PYG{o}{=} \PYG{l+s+s2}{\PYGZdq{}}\PYG{l+s+s2}{/home/sadi/LSST/test/wlcf\PYGZhy{}main/tests/Bell\PYGZus{}outputs}\PYG{l+s+s2}{\PYGZdq{}}
\PYG{n}{prefix} \PYG{o}{=} \PYG{l+s+s2}{\PYGZdq{}}\PYG{l+s+s2}{halo\PYGZus{}model\PYGZus{}z05078\PYGZus{}}\PYG{l+s+s2}{\PYGZdq{}}

\PYG{n}{theta} \PYG{o}{=} \PYG{n}{np}\PYG{o}{.}\PYG{n}{loadtxt}\PYG{p}{(}\PYG{l+s+sa}{f}\PYG{l+s+s2}{\PYGZdq{}}\PYG{l+s+si}{\PYGZob{}}\PYG{n}{path}\PYG{l+s+si}{\PYGZcb{}}\PYG{l+s+s2}{/}\PYG{l+s+si}{\PYGZob{}}\PYG{n}{prefix}\PYG{l+s+si}{\PYGZcb{}}\PYG{l+s+s2}{theta\PYGZus{}array.txt}\PYG{l+s+s2}{\PYGZdq{}}\PYG{p}{)} \PYG{o}{*} \PYG{l+m+mi}{180}\PYG{o}{/}\PYG{n}{np}\PYG{o}{.}\PYG{n}{pi} \PYG{o}{*} \PYG{l+m+mi}{60}

\PYG{n}{theta2\PYGZus{}targets} \PYG{o}{=} \PYG{p}{[}\PYG{l+m+mi}{13}\PYG{p}{,} \PYG{l+m+mi}{72}\PYG{p}{,} \PYG{l+m+mi}{169}\PYG{p}{]}
\PYG{n}{m\PYGZus{}values} \PYG{o}{=} \PYG{p}{[}\PYG{l+m+mi}{0}\PYG{p}{,} \PYG{l+m+mi}{1}\PYG{p}{,} \PYG{l+m+mi}{2}\PYG{p}{]}

\PYG{n}{fig}\PYG{p}{,} \PYG{n}{axes} \PYG{o}{=} \PYG{n}{plt}\PYG{o}{.}\PYG{n}{subplots}\PYG{p}{(}
    \PYG{n+nb}{len}\PYG{p}{(}\PYG{n}{theta2\PYGZus{}targets}\PYG{p}{)}\PYG{p}{,}
    \PYG{n+nb}{len}\PYG{p}{(}\PYG{n}{m\PYGZus{}values}\PYG{p}{)}\PYG{p}{,}
    \PYG{n}{figsize}\PYG{o}{=}\PYG{p}{(}\PYG{l+m+mi}{9}\PYG{p}{,} \PYG{l+m+mi}{9}\PYG{p}{)}\PYG{p}{,}
    \PYG{n}{sharex}\PYG{o}{=}\PYG{k+kc}{True}
\PYG{p}{)}

\PYG{k}{for} \PYG{n}{i}\PYG{p}{,} \PYG{n}{theta2\PYGZus{}fixed} \PYG{o+ow}{in} \PYG{n+nb}{enumerate}\PYG{p}{(}\PYG{n}{theta2\PYGZus{}targets}\PYG{p}{)}\PYG{p}{:}
    \PYG{n}{idx} \PYG{o}{=} \PYG{n}{np}\PYG{o}{.}\PYG{n}{argmin}\PYG{p}{(}\PYG{n}{np}\PYG{o}{.}\PYG{n}{abs}\PYG{p}{(}\PYG{n}{theta} \PYG{o}{\PYGZhy{}} \PYG{n}{theta2\PYGZus{}fixed}\PYG{p}{)}\PYG{p}{)}
    \PYG{n}{theta2\PYGZus{}actual} \PYG{o}{=} \PYG{n}{theta}\PYG{p}{[}\PYG{n}{idx}\PYG{p}{]}

    \PYG{k}{for} \PYG{n}{j}\PYG{p}{,} \PYG{n}{m} \PYG{o+ow}{in} \PYG{n+nb}{enumerate}\PYG{p}{(}\PYG{n}{m\PYGZus{}values}\PYG{p}{)}\PYG{p}{:}
        \PYG{n}{ax} \PYG{o}{=} \PYG{n}{axes}\PYG{p}{[}\PYG{n}{i}\PYG{p}{,} \PYG{n}{j}\PYG{p}{]}

        \PYG{n}{zeta} \PYG{o}{=} \PYG{n}{np}\PYG{o}{.}\PYG{n}{abs}\PYG{p}{(}\PYG{n}{np}\PYG{o}{.}\PYG{n}{loadtxt}\PYG{p}{(}\PYG{l+s+sa}{f}\PYG{l+s+s2}{\PYGZdq{}}\PYG{l+s+si}{\PYGZob{}}\PYG{n}{path}\PYG{l+s+si}{\PYGZcb{}}\PYG{l+s+s2}{/}\PYG{l+s+si}{\PYGZob{}}\PYG{n}{prefix}\PYG{l+s+si}{\PYGZcb{}}\PYG{l+s+s2}{zetam}\PYG{l+s+si}{\PYGZob{}}\PYG{n}{m}\PYG{l+s+si}{\PYGZcb{}}\PYG{l+s+s2}{.txt}\PYG{l+s+s2}{\PYGZdq{}}\PYG{p}{)}\PYG{p}{)}

        \PYG{c+c1}{\PYGZsh{} theta2 fixed, theta1 varies}
        \PYG{n}{y} \PYG{o}{=} \PYG{n}{zeta}\PYG{p}{[}\PYG{p}{:}\PYG{p}{,} \PYG{n}{idx}\PYG{p}{]}

        \PYG{n}{ax}\PYG{o}{.}\PYG{n}{plot}\PYG{p}{(}\PYG{n}{theta}\PYG{p}{,} \PYG{n}{y}\PYG{p}{,} \PYG{n}{color}\PYG{o}{=}\PYG{l+s+s2}{\PYGZdq{}}\PYG{l+s+s2}{black}\PYG{l+s+s2}{\PYGZdq{}}\PYG{p}{,} \PYG{n}{marker}\PYG{o}{=}\PYG{l+s+s2}{\PYGZdq{}}\PYG{l+s+s2}{o}\PYG{l+s+s2}{\PYGZdq{}}\PYG{p}{,} \PYG{n}{markersize}\PYG{o}{=}\PYG{l+m+mi}{3}\PYG{p}{,} \PYG{n}{lw}\PYG{o}{=}\PYG{l+m+mi}{1}\PYG{p}{)}

        \PYG{n}{ax}\PYG{o}{.}\PYG{n}{set\PYGZus{}xscale}\PYG{p}{(}\PYG{l+s+s2}{\PYGZdq{}}\PYG{l+s+s2}{log}\PYG{l+s+s2}{\PYGZdq{}}\PYG{p}{)}
        \PYG{n}{ax}\PYG{o}{.}\PYG{n}{set\PYGZus{}yscale}\PYG{p}{(}\PYG{l+s+s2}{\PYGZdq{}}\PYG{l+s+s2}{log}\PYG{l+s+s2}{\PYGZdq{}}\PYG{p}{)}

        \PYG{n}{ax}\PYG{o}{.}\PYG{n}{set\PYGZus{}xlim}\PYG{p}{(}\PYG{l+m+mi}{10}\PYG{p}{,} \PYG{l+m+mi}{200}\PYG{p}{)}

        \PYG{k}{if} \PYG{n}{i} \PYG{o}{==} \PYG{l+m+mi}{0}\PYG{p}{:}
            \PYG{n}{ax}\PYG{o}{.}\PYG{n}{set\PYGZus{}title}\PYG{p}{(}\PYG{l+s+sa}{fr}\PYG{l+s+s2}{\PYGZdq{}}\PYG{l+s+s2}{\PYGZdl{}m=}\PYG{l+s+si}{\PYGZob{}}\PYG{n}{m}\PYG{l+s+si}{\PYGZcb{}}\PYG{l+s+s2}{\PYGZdl{}}\PYG{l+s+s2}{\PYGZdq{}}\PYG{p}{)}

        \PYG{k}{if} \PYG{n}{j} \PYG{o}{==} \PYG{l+m+mi}{0}\PYG{p}{:}
            \PYG{n}{ax}\PYG{o}{.}\PYG{n}{set\PYGZus{}ylabel}\PYG{p}{(}
                \PYG{l+s+sa}{fr}\PYG{l+s+s2}{\PYGZdq{}}\PYG{l+s+s2}{\PYGZdl{}}\PYG{l+s+s2}{\PYGZbs{}}\PYG{l+s+s2}{zeta\PYGZus{}m(}\PYG{l+s+s2}{\PYGZbs{}}\PYG{l+s+s2}{theta\PYGZus{}1,}\PYG{l+s+s2}{\PYGZbs{}}\PYG{l+s+s2}{theta\PYGZus{}2=}\PYG{l+s+si}{\PYGZob{}}\PYG{n}{theta2\PYGZus{}actual}\PYG{l+s+si}{:}\PYG{l+s+s2}{.0f}\PYG{l+s+si}{\PYGZcb{}}\PYG{l+s+s2}{\PYGZsq{}}\PYG{l+s+s2}{)\PYGZdl{}}\PYG{l+s+s2}{\PYGZdq{}}
            \PYG{p}{)}

        \PYG{k}{if} \PYG{n}{i} \PYG{o}{==} \PYG{n+nb}{len}\PYG{p}{(}\PYG{n}{theta2\PYGZus{}targets}\PYG{p}{)} \PYG{o}{\PYGZhy{}} \PYG{l+m+mi}{1}\PYG{p}{:}
            \PYG{n}{ax}\PYG{o}{.}\PYG{n}{set\PYGZus{}xlabel}\PYG{p}{(}\PYG{l+s+sa}{r}\PYG{l+s+s2}{\PYGZdq{}}\PYG{l+s+s2}{\PYGZdl{}}\PYG{l+s+s2}{\PYGZbs{}}\PYG{l+s+s2}{theta\PYGZus{}1\PYGZdl{} [arcmin]}\PYG{l+s+s2}{\PYGZdq{}}\PYG{p}{)}

        \PYG{n}{ax}\PYG{o}{.}\PYG{n}{grid}\PYG{p}{(}\PYG{n}{alpha}\PYG{o}{=}\PYG{l+m+mf}{0.2}\PYG{p}{)}

\PYG{n}{plt}\PYG{o}{.}\PYG{n}{tight\PYGZus{}layout}\PYG{p}{(}\PYG{p}{)}
\PYG{n}{plt}\PYG{o}{.}\PYG{n}{savefig}\PYG{p}{(}\PYG{l+s+s2}{\PYGZdq{}}\PYG{l+s+s2}{halo\PYGZus{}model\PYGZus{}slice\PYGZus{}3pcf.png}\PYG{l+s+s2}{\PYGZdq{}}\PYG{p}{,} \PYG{n}{dpi}\PYG{o}{=}\PYG{l+m+mi}{300}\PYG{p}{,} \PYG{n}{bbox\PYGZus{}inches}\PYG{o}{=}\PYG{l+s+s2}{\PYGZdq{}}\PYG{l+s+s2}{tight}\PYG{l+s+s2}{\PYGZdq{}}\PYG{p}{)}
\PYG{n}{plt}\PYG{o}{.}\PYG{n}{show}\PYG{p}{(}\PYG{p}{)}
\end{sphinxVerbatim}

\noindent{\hspace*{\fill}\sphinxincludegraphics[width=0.700\linewidth]{{halo_model_slice_3pcf}.png}\hspace*{\fill}}


\section{Comparing different models}
\label{\detokenize{handson:comparing-different-models}}
\sphinxAtStartPar
We can run several theoretical models by changing
\sphinxcode{\sphinxupquote{tree\_level}}:

\begin{sphinxVerbatim}[commandchars=\\\{\}]
\PYG{n}{base\PYGZus{}params} \PYG{o}{=} \PYG{n+nb}{dict}\PYG{p}{(}
    \PYG{n}{numberThreads}\PYG{o}{=}\PYG{l+m+mi}{8}\PYG{p}{,}
    \PYG{n}{verbose}\PYG{o}{=}\PYG{l+m+mi}{2}\PYG{p}{,}
    \PYG{n}{verbose\PYGZus{}log}\PYG{o}{=}\PYG{l+m+mi}{1}\PYG{p}{,}

    \PYG{c+c1}{\PYGZsh{} T17 / paper cosmology}
    \PYG{n}{z}\PYG{o}{=}\PYG{l+m+mf}{0.5078}\PYG{p}{,}
    \PYG{n}{h}\PYG{o}{=}\PYG{l+m+mf}{0.7}\PYG{p}{,}
    \PYG{n}{sigma8}\PYG{o}{=}\PYG{l+m+mf}{0.82}\PYG{p}{,}
    \PYG{n}{Omb}\PYG{o}{=}\PYG{l+m+mf}{0.046}\PYG{p}{,}
    \PYG{n}{Omc}\PYG{o}{=}\PYG{l+m+mf}{0.233}\PYG{p}{,}
    \PYG{n}{Omnu}\PYG{o}{=}\PYG{l+m+mf}{0.0}\PYG{p}{,}
    \PYG{n}{ns}\PYG{o}{=}\PYG{l+m+mf}{0.97}\PYG{p}{,}
    \PYG{n}{w}\PYG{o}{=}\PYG{o}{\PYGZhy{}}\PYG{l+m+mf}{1.0}\PYG{p}{,}

    \PYG{c+c1}{\PYGZsh{} Make sure this P(k) was generated with the same cosmology}
    \PYG{n}{fnamePS}\PYG{o}{=}\PYG{l+s+s2}{\PYGZdq{}}\PYG{l+s+s2}{./input/linear\PYGZus{}pk\PYGZus{}camb\PYGZus{}T17\PYGZus{}z05078.txt}\PYG{l+s+s2}{\PYGZdq{}}\PYG{p}{,}

    \PYG{c+c1}{\PYGZsh{} Redshift bin}
    \PYG{n}{Wg}\PYG{o}{=}\PYG{l+m+mi}{0}\PYG{p}{,}
    \PYG{n}{fWgchi}\PYG{o}{=}\PYG{l+s+s2}{\PYGZdq{}}\PYG{l+s+s2}{./input/Wg\PYGZus{}Takahashi\PYGZus{}z05078.txt}\PYG{l+s+s2}{\PYGZdq{}}\PYG{p}{,}

    \PYG{n}{rootDir}\PYG{o}{=}\PYG{l+s+s2}{\PYGZdq{}}\PYG{l+s+s2}{Output}\PYG{l+s+s2}{\PYGZdq{}}\PYG{p}{,}
    \PYG{n}{path\PYGZus{}Bells}\PYG{o}{=}\PYG{l+s+s2}{\PYGZdq{}}\PYG{l+s+s2}{Bell\PYGZus{}outputs}\PYG{l+s+s2}{\PYGZdq{}}\PYG{p}{,}

    \PYG{n}{zbin}\PYG{o}{=}\PYG{l+m+mf}{0.5078}\PYG{p}{,}
    \PYG{n}{mMax}\PYG{o}{=}\PYG{l+m+mi}{5}\PYG{p}{,}
    \PYG{n}{chiQuadSteps}\PYG{o}{=}\PYG{l+m+mi}{300}\PYG{p}{,}
    \PYG{n}{GLpoints}\PYG{o}{=}\PYG{l+m+mi}{64}\PYG{p}{,}
    \PYG{n}{Nell}\PYG{o}{=}\PYG{l+m+mi}{128}\PYG{p}{,}
    \PYG{n}{ellmax}\PYG{o}{=}\PYG{l+m+mi}{10000}\PYG{p}{,}
    \PYG{n}{ellmin}\PYG{o}{=}\PYG{l+m+mf}{0.001}\PYG{p}{,}
    \PYG{n}{writevectors}\PYG{o}{=}\PYG{l+m+mi}{0}\PYG{p}{,}
    \PYG{n}{chatty}\PYG{o}{=}\PYG{l+m+mi}{2}\PYG{p}{,}
    \PYG{n}{options}\PYG{o}{=}\PYG{l+s+s2}{\PYGZdq{}}\PYG{l+s+s2}{\PYGZdq{}}
\PYG{p}{)}

\PYG{n}{models} \PYG{o}{=} \PYG{p}{\PYGZob{}}
    \PYG{l+m+mi}{1}\PYG{p}{:} \PYG{l+s+s2}{\PYGZdq{}}\PYG{l+s+s2}{SPT}\PYG{l+s+s2}{\PYGZdq{}}\PYG{p}{,}
    \PYG{l+m+mi}{2}\PYG{p}{:} \PYG{l+s+s2}{\PYGZdq{}}\PYG{l+s+s2}{Tree}\PYG{l+s+s2}{\PYGZdq{}}\PYG{p}{,}
    \PYG{l+m+mi}{3}\PYG{p}{:} \PYG{l+s+s2}{\PYGZdq{}}\PYG{l+s+s2}{EFT}\PYG{l+s+s2}{\PYGZdq{}}\PYG{p}{,}
    \PYG{l+m+mi}{4}\PYG{p}{:} \PYG{l+s+s2}{\PYGZdq{}}\PYG{l+s+s2}{Halo\PYGZus{}model}\PYG{l+s+s2}{\PYGZdq{}}
\PYG{p}{\PYGZcb{}}

\PYG{k}{for} \PYG{n}{tree\PYGZus{}level}\PYG{p}{,} \PYG{n}{model\PYGZus{}name} \PYG{o+ow}{in} \PYG{n}{models}\PYG{o}{.}\PYG{n}{items}\PYG{p}{(}\PYG{p}{)}\PYG{p}{:}

    \PYG{n+nb}{print}\PYG{p}{(}\PYG{l+s+sa}{f}\PYG{l+s+s2}{\PYGZdq{}}\PYG{l+s+s2}{Running }\PYG{l+s+si}{\PYGZob{}}\PYG{n}{model\PYGZus{}name}\PYG{l+s+si}{\PYGZcb{}}\PYG{l+s+s2}{ model with tree\PYGZus{}level=}\PYG{l+s+si}{\PYGZob{}}\PYG{n}{tree\PYGZus{}level}\PYG{l+s+si}{\PYGZcb{}}\PYG{l+s+s2}{\PYGZdq{}}\PYG{p}{)}

    \PYG{n}{params} \PYG{o}{=} \PYG{n}{base\PYGZus{}params}\PYG{o}{.}\PYG{n}{copy}\PYG{p}{(}\PYG{p}{)}
    \PYG{n}{params}\PYG{p}{[}\PYG{l+s+s2}{\PYGZdq{}}\PYG{l+s+s2}{tree\PYGZus{}level}\PYG{l+s+s2}{\PYGZdq{}}\PYG{p}{]} \PYG{o}{=} \PYG{n}{tree\PYGZus{}level}
    \PYG{n}{params}\PYG{p}{[}\PYG{l+s+s2}{\PYGZdq{}}\PYG{l+s+s2}{prefix}\PYG{l+s+s2}{\PYGZdq{}}\PYG{p}{]} \PYG{o}{=} \PYG{l+s+sa}{f}\PYG{l+s+s2}{\PYGZdq{}}\PYG{l+s+si}{\PYGZob{}}\PYG{n}{model\PYGZus{}name}\PYG{o}{.}\PYG{n}{lower}\PYG{p}{(}\PYG{p}{)}\PYG{l+s+si}{\PYGZcb{}}\PYG{l+s+s2}{\PYGZus{}z05078\PYGZus{}}\PYG{l+s+s2}{\PYGZdq{}}

    \PYG{n}{w} \PYG{o}{=} \PYG{n}{wlcf}\PYG{p}{(}\PYG{p}{)}
    \PYG{n}{w}\PYG{o}{.}\PYG{n}{set}\PYG{p}{(}\PYG{n}{params}\PYG{p}{)}
    \PYG{n}{w}\PYG{o}{.}\PYG{n}{Run}\PYG{p}{(}\PYG{p}{)}

    \PYG{n+nb}{print}\PYG{p}{(}\PYG{l+s+sa}{f}\PYG{l+s+s2}{\PYGZdq{}}\PYG{l+s+s2}{Finished }\PYG{l+s+si}{\PYGZob{}}\PYG{n}{model\PYGZus{}name}\PYG{l+s+si}{\PYGZcb{}}\PYG{l+s+s2}{\PYGZdq{}}\PYG{p}{)}
\end{sphinxVerbatim}

\sphinxAtStartPar
We plot the different cases

\begin{sphinxVerbatim}[commandchars=\\\{\}]
\PYG{n}{plt}\PYG{o}{.}\PYG{n}{rcParams}\PYG{o}{.}\PYG{n}{update}\PYG{p}{(}\PYG{p}{\PYGZob{}}
    \PYG{l+s+s2}{\PYGZdq{}}\PYG{l+s+s2}{font.size}\PYG{l+s+s2}{\PYGZdq{}}\PYG{p}{:} \PYG{l+m+mi}{14}\PYG{p}{,}
    \PYG{l+s+s2}{\PYGZdq{}}\PYG{l+s+s2}{axes.titlesize}\PYG{l+s+s2}{\PYGZdq{}}\PYG{p}{:} \PYG{l+m+mi}{16}\PYG{p}{,}
    \PYG{l+s+s2}{\PYGZdq{}}\PYG{l+s+s2}{axes.labelsize}\PYG{l+s+s2}{\PYGZdq{}}\PYG{p}{:} \PYG{l+m+mi}{14}
\PYG{p}{\PYGZcb{}}\PYG{p}{)}

\PYG{n}{path} \PYG{o}{=} \PYG{l+s+s2}{\PYGZdq{}}\PYG{l+s+s2}{/home/sadi/LSST/test/wlcf\PYGZhy{}main/tests/Bell\PYGZus{}outputs}\PYG{l+s+s2}{\PYGZdq{}}

\PYG{n}{models} \PYG{o}{=} \PYG{p}{\PYGZob{}}
    \PYG{l+s+s2}{\PYGZdq{}}\PYG{l+s+s2}{SPT}\PYG{l+s+s2}{\PYGZdq{}}\PYG{p}{:} \PYG{l+s+s2}{\PYGZdq{}}\PYG{l+s+s2}{spt\PYGZus{}z05078\PYGZus{}}\PYG{l+s+s2}{\PYGZdq{}}\PYG{p}{,}
    \PYG{l+s+s2}{\PYGZdq{}}\PYG{l+s+s2}{Tree}\PYG{l+s+s2}{\PYGZdq{}}\PYG{p}{:} \PYG{l+s+s2}{\PYGZdq{}}\PYG{l+s+s2}{tree\PYGZus{}z05078\PYGZus{}}\PYG{l+s+s2}{\PYGZdq{}}\PYG{p}{,}
    \PYG{l+s+s2}{\PYGZdq{}}\PYG{l+s+s2}{EFT}\PYG{l+s+s2}{\PYGZdq{}}\PYG{p}{:} \PYG{l+s+s2}{\PYGZdq{}}\PYG{l+s+s2}{eft\PYGZus{}z05078\PYGZus{}}\PYG{l+s+s2}{\PYGZdq{}}\PYG{p}{,}
    \PYG{l+s+s2}{\PYGZdq{}}\PYG{l+s+s2}{Halo Model}\PYG{l+s+s2}{\PYGZdq{}}\PYG{p}{:} \PYG{l+s+s2}{\PYGZdq{}}\PYG{l+s+s2}{halo\PYGZus{}model\PYGZus{}z05078\PYGZus{}}\PYG{l+s+s2}{\PYGZdq{}}\PYG{p}{,}
\PYG{p}{\PYGZcb{}}

\PYG{n}{m\PYGZus{}values} \PYG{o}{=} \PYG{n+nb}{range}\PYG{p}{(}\PYG{l+m+mi}{5}\PYG{p}{)}

\PYG{n}{theta} \PYG{o}{=} \PYG{n}{np}\PYG{o}{.}\PYG{n}{loadtxt}\PYG{p}{(}\PYG{l+s+sa}{f}\PYG{l+s+s2}{\PYGZdq{}}\PYG{l+s+si}{\PYGZob{}}\PYG{n}{path}\PYG{l+s+si}{\PYGZcb{}}\PYG{l+s+s2}{/spt\PYGZus{}z05\PYGZus{}theta\PYGZus{}array.txt}\PYG{l+s+s2}{\PYGZdq{}}\PYG{p}{)} \PYG{o}{*} \PYG{l+m+mi}{180} \PYG{o}{/} \PYG{n}{np}\PYG{o}{.}\PYG{n}{pi} \PYG{o}{*} \PYG{l+m+mi}{60}
\PYG{n}{Theta2}\PYG{p}{,} \PYG{n}{Theta1} \PYG{o}{=} \PYG{n}{np}\PYG{o}{.}\PYG{n}{meshgrid}\PYG{p}{(}\PYG{n}{theta}\PYG{p}{,} \PYG{n}{theta}\PYG{p}{)}

\PYG{n}{nrows} \PYG{o}{=} \PYG{n+nb}{len}\PYG{p}{(}\PYG{n}{models}\PYG{p}{)}
\PYG{n}{ncols} \PYG{o}{=} \PYG{n+nb}{len}\PYG{p}{(}\PYG{n}{m\PYGZus{}values}\PYG{p}{)}

\PYG{n}{fig} \PYG{o}{=} \PYG{n}{plt}\PYG{o}{.}\PYG{n}{figure}\PYG{p}{(}\PYG{n}{figsize}\PYG{o}{=}\PYG{p}{(}\PYG{l+m+mi}{20}\PYG{p}{,} \PYG{l+m+mi}{13}\PYG{p}{)}\PYG{p}{)}

\PYG{n}{gs} \PYG{o}{=} \PYG{n}{fig}\PYG{o}{.}\PYG{n}{add\PYGZus{}gridspec}\PYG{p}{(}
    \PYG{n}{nrows}\PYG{o}{=}\PYG{n}{nrows}\PYG{p}{,}
    \PYG{n}{ncols}\PYG{o}{=}\PYG{n}{ncols} \PYG{o}{+} \PYG{l+m+mi}{1}\PYG{p}{,}
    \PYG{n}{width\PYGZus{}ratios}\PYG{o}{=}\PYG{p}{[}\PYG{l+m+mi}{1}\PYG{p}{,} \PYG{l+m+mi}{1}\PYG{p}{,} \PYG{l+m+mi}{1}\PYG{p}{,} \PYG{l+m+mi}{1}\PYG{p}{,} \PYG{l+m+mi}{1}\PYG{p}{,} \PYG{l+m+mf}{0.05}\PYG{p}{]}\PYG{p}{,}
    \PYG{n}{wspace}\PYG{o}{=}\PYG{l+m+mf}{0.2}\PYG{p}{,}
    \PYG{n}{hspace}\PYG{o}{=}\PYG{l+m+mf}{0.25}
\PYG{p}{)}

\PYG{n}{norm} \PYG{o}{=} \PYG{n}{LogNorm}\PYG{p}{(}\PYG{n}{vmin}\PYG{o}{=}\PYG{l+m+mf}{1e\PYGZhy{}11}\PYG{p}{,} \PYG{n}{vmax}\PYG{o}{=}\PYG{l+m+mf}{1e\PYGZhy{}8}\PYG{p}{)}

\PYG{k}{for} \PYG{n}{i}\PYG{p}{,} \PYG{p}{(}\PYG{n}{model\PYGZus{}name}\PYG{p}{,} \PYG{n}{prefix}\PYG{p}{)} \PYG{o+ow}{in} \PYG{n+nb}{enumerate}\PYG{p}{(}\PYG{n}{models}\PYG{o}{.}\PYG{n}{items}\PYG{p}{(}\PYG{p}{)}\PYG{p}{)}\PYG{p}{:}
    \PYG{k}{for} \PYG{n}{j}\PYG{p}{,} \PYG{n}{m} \PYG{o+ow}{in} \PYG{n+nb}{enumerate}\PYG{p}{(}\PYG{n}{m\PYGZus{}values}\PYG{p}{)}\PYG{p}{:}

        \PYG{n}{ax} \PYG{o}{=} \PYG{n}{fig}\PYG{o}{.}\PYG{n}{add\PYGZus{}subplot}\PYG{p}{(}\PYG{n}{gs}\PYG{p}{[}\PYG{n}{i}\PYG{p}{,} \PYG{n}{j}\PYG{p}{]}\PYG{p}{)}

        \PYG{n}{zeta} \PYG{o}{=} \PYG{n}{np}\PYG{o}{.}\PYG{n}{abs}\PYG{p}{(}\PYG{n}{np}\PYG{o}{.}\PYG{n}{loadtxt}\PYG{p}{(}\PYG{l+s+sa}{f}\PYG{l+s+s2}{\PYGZdq{}}\PYG{l+s+si}{\PYGZob{}}\PYG{n}{path}\PYG{l+s+si}{\PYGZcb{}}\PYG{l+s+s2}{/}\PYG{l+s+si}{\PYGZob{}}\PYG{n}{prefix}\PYG{l+s+si}{\PYGZcb{}}\PYG{l+s+s2}{zetam}\PYG{l+s+si}{\PYGZob{}}\PYG{n}{m}\PYG{l+s+si}{\PYGZcb{}}\PYG{l+s+s2}{.txt}\PYG{l+s+s2}{\PYGZdq{}}\PYG{p}{)}\PYG{p}{)}

        \PYG{n}{im} \PYG{o}{=} \PYG{n}{ax}\PYG{o}{.}\PYG{n}{pcolormesh}\PYG{p}{(}
            \PYG{n}{Theta2}\PYG{p}{,}
            \PYG{n}{Theta1}\PYG{p}{,}
            \PYG{n}{zeta}\PYG{p}{,}
            \PYG{n}{shading}\PYG{o}{=}\PYG{l+s+s2}{\PYGZdq{}}\PYG{l+s+s2}{auto}\PYG{l+s+s2}{\PYGZdq{}}\PYG{p}{,}
            \PYG{n}{cmap}\PYG{o}{=}\PYG{l+s+s2}{\PYGZdq{}}\PYG{l+s+s2}{RdYlBu\PYGZus{}r}\PYG{l+s+s2}{\PYGZdq{}}\PYG{p}{,}
            \PYG{n}{norm}\PYG{o}{=}\PYG{n}{norm}
        \PYG{p}{)}

        \PYG{n}{ax}\PYG{o}{.}\PYG{n}{set\PYGZus{}xscale}\PYG{p}{(}\PYG{l+s+s2}{\PYGZdq{}}\PYG{l+s+s2}{log}\PYG{l+s+s2}{\PYGZdq{}}\PYG{p}{)}
        \PYG{n}{ax}\PYG{o}{.}\PYG{n}{set\PYGZus{}yscale}\PYG{p}{(}\PYG{l+s+s2}{\PYGZdq{}}\PYG{l+s+s2}{log}\PYG{l+s+s2}{\PYGZdq{}}\PYG{p}{)}
        \PYG{n}{ax}\PYG{o}{.}\PYG{n}{set\PYGZus{}xlim}\PYG{p}{(}\PYG{l+m+mi}{10}\PYG{p}{,} \PYG{l+m+mi}{200}\PYG{p}{)}
        \PYG{n}{ax}\PYG{o}{.}\PYG{n}{set\PYGZus{}ylim}\PYG{p}{(}\PYG{l+m+mi}{10}\PYG{p}{,} \PYG{l+m+mi}{200}\PYG{p}{)}

        \PYG{k}{if} \PYG{n}{i} \PYG{o}{==} \PYG{l+m+mi}{0}\PYG{p}{:}
            \PYG{n}{ax}\PYG{o}{.}\PYG{n}{set\PYGZus{}title}\PYG{p}{(}\PYG{l+s+sa}{fr}\PYG{l+s+s2}{\PYGZdq{}}\PYG{l+s+s2}{\PYGZdl{}m=}\PYG{l+s+si}{\PYGZob{}}\PYG{n}{m}\PYG{l+s+si}{\PYGZcb{}}\PYG{l+s+s2}{\PYGZdl{}}\PYG{l+s+s2}{\PYGZdq{}}\PYG{p}{)}

        \PYG{k}{if} \PYG{n}{j} \PYG{o}{==} \PYG{l+m+mi}{0}\PYG{p}{:}
            \PYG{n}{ax}\PYG{o}{.}\PYG{n}{set\PYGZus{}ylabel}\PYG{p}{(}\PYG{l+s+sa}{f}\PYG{l+s+s2}{\PYGZdq{}}\PYG{l+s+si}{\PYGZob{}}\PYG{n}{model\PYGZus{}name}\PYG{l+s+si}{\PYGZcb{}}\PYG{l+s+se}{\PYGZbs{}n}\PYG{l+s+s2}{\PYGZdq{}} \PYG{o}{+} \PYG{l+s+sa}{r}\PYG{l+s+s2}{\PYGZdq{}}\PYG{l+s+s2}{\PYGZdl{}}\PYG{l+s+s2}{\PYGZbs{}}\PYG{l+s+s2}{theta\PYGZus{}1\PYGZdl{} [arcmin]}\PYG{l+s+s2}{\PYGZdq{}}\PYG{p}{)}

        \PYG{k}{if} \PYG{n}{i} \PYG{o}{==} \PYG{n}{nrows} \PYG{o}{\PYGZhy{}} \PYG{l+m+mi}{1}\PYG{p}{:}
            \PYG{n}{ax}\PYG{o}{.}\PYG{n}{set\PYGZus{}xlabel}\PYG{p}{(}\PYG{l+s+sa}{r}\PYG{l+s+s2}{\PYGZdq{}}\PYG{l+s+s2}{\PYGZdl{}}\PYG{l+s+s2}{\PYGZbs{}}\PYG{l+s+s2}{theta\PYGZus{}2\PYGZdl{} [arcmin]}\PYG{l+s+s2}{\PYGZdq{}}\PYG{p}{)}

\PYG{n}{cax} \PYG{o}{=} \PYG{n}{fig}\PYG{o}{.}\PYG{n}{add\PYGZus{}subplot}\PYG{p}{(}\PYG{n}{gs}\PYG{p}{[}\PYG{p}{:}\PYG{p}{,} \PYG{o}{\PYGZhy{}}\PYG{l+m+mi}{1}\PYG{p}{]}\PYG{p}{)}
\PYG{n}{cbar} \PYG{o}{=} \PYG{n}{fig}\PYG{o}{.}\PYG{n}{colorbar}\PYG{p}{(}\PYG{n}{im}\PYG{p}{,} \PYG{n}{cax}\PYG{o}{=}\PYG{n}{cax}\PYG{p}{)}
\PYG{n}{cbar}\PYG{o}{.}\PYG{n}{set\PYGZus{}label}\PYG{p}{(}\PYG{l+s+sa}{r}\PYG{l+s+s2}{\PYGZdq{}}\PYG{l+s+s2}{\PYGZdl{}|}\PYG{l+s+s2}{\PYGZbs{}}\PYG{l+s+s2}{zeta\PYGZus{}m|\PYGZdl{}}\PYG{l+s+s2}{\PYGZdq{}}\PYG{p}{,} \PYG{n}{fontsize}\PYG{o}{=}\PYG{l+m+mi}{16}\PYG{p}{)}
\PYG{n}{cbar}\PYG{o}{.}\PYG{n}{ax}\PYG{o}{.}\PYG{n}{tick\PYGZus{}params}\PYG{p}{(}\PYG{n}{labelsize}\PYG{o}{=}\PYG{l+m+mi}{12}\PYG{p}{)}
\PYG{n}{plt}\PYG{o}{.}\PYG{n}{savefig}\PYG{p}{(}\PYG{l+s+s2}{\PYGZdq{}}\PYG{l+s+s2}{models\PYGZus{}3pcf.png}\PYG{l+s+s2}{\PYGZdq{}}\PYG{p}{,} \PYG{n}{dpi}\PYG{o}{=}\PYG{l+m+mi}{300}\PYG{p}{,} \PYG{n}{bbox\PYGZus{}inches}\PYG{o}{=}\PYG{l+s+s2}{\PYGZdq{}}\PYG{l+s+s2}{tight}\PYG{l+s+s2}{\PYGZdq{}}\PYG{p}{)}
\PYG{n}{plt}\PYG{o}{.}\PYG{n}{show}\PYG{p}{(}\PYG{p}{)}
\end{sphinxVerbatim}

\noindent{\hspace*{\fill}\sphinxincludegraphics[width=1.000\linewidth]{{models_3pcf}.png}\hspace*{\fill}}

\sphinxstepscope


\chapter{3\sphinxhyphen{}point correlation functions}
\label{\detokenize{3pcf:point-correlation-functions}}\label{\detokenize{3pcf::doc}}
\sphinxAtStartPar
\sphinxstylestrong{wlcf} computes \sphinxstylestrong{three\sphinxhyphen{}point correlation functions (3PCF)} for scalar fields, such as the weak lensing convergence \(\kappa\). These higher\sphinxhyphen{}order statistics provide access to non\sphinxhyphen{}Gaussian information in the large\sphinxhyphen{}scale structure of the Universe.

\sphinxAtStartPar
In particular, the code evaluates the correlation:
\begin{equation*}
\begin{split}\langle \kappa(\hat{n}_1),\kappa(\hat{n}_2),\kappa(\hat{n}_3) \rangle\end{split}
\end{equation*}
\sphinxAtStartPar
using a plane\sphinxhyphen{}wave expansion and a harmonic decomposition.

\sphinxAtStartPar
\sphinxstylestrong{Harmonic base}

\sphinxAtStartPar
The angular dependence of the 3PCF is expressed in a \sphinxstylestrong{harmonic basis}, where the signal is decomposed into multipoles. This allows for an efficient and structured representation of the correlation function.

\sphinxAtStartPar
The expansion is performed in terms of angular modes labeled by an index \(m\), leading to a set of coefficients:
\begin{equation*}
\begin{split}\zeta_m(\theta_1, \theta_2)\end{split}
\end{equation*}
\sphinxAtStartPar
where:
\begin{itemize}
\item {} 
\sphinxAtStartPar
\(\theta_1\), \(\theta_2\) are angular separations

\item {} 
\sphinxAtStartPar
\(m = 0, 1, \dots, m_{\mathrm{max}}\) is the multipole index

\end{itemize}

\sphinxAtStartPar
\sphinxstylestrong{Output}

\sphinxAtStartPar
The results of the harmonic decomposition are stored in files:

\begin{sphinxVerbatim}[commandchars=\\\{\}]
\PYG{n}{zetam0}\PYG{o}{.}\PYG{n}{txt}
\PYG{n}{zetam1}\PYG{o}{.}\PYG{n}{txt}
\PYG{o}{.}\PYG{o}{.}\PYG{o}{.}
\PYG{n}{zetamN}\PYG{o}{.}\PYG{n}{txt}
\end{sphinxVerbatim}

\sphinxAtStartPar
along with the corresponding angular grid:

\begin{sphinxVerbatim}[commandchars=\\\{\}]
\PYG{n}{theta\PYGZus{}array}\PYG{o}{.}\PYG{n}{txt}
\end{sphinxVerbatim}

\sphinxAtStartPar
These outputs can be used to reconstruct or visualize the full three\sphinxhyphen{}point correlation function.


\chapter{Indices and tables}
\label{\detokenize{index:indices-and-tables}}\begin{itemize}
\item {} 
\sphinxAtStartPar
\DUrole{xref}{\DUrole{std}{\DUrole{std-ref}{genindex}}}

\item {} 
\sphinxAtStartPar
\DUrole{xref}{\DUrole{std}{\DUrole{std-ref}{modindex}}}

\item {} 
\sphinxAtStartPar
\DUrole{xref}{\DUrole{std}{\DUrole{std-ref}{search}}}

\end{itemize}



\renewcommand{\indexname}{Index}
\printindex
\end{document}